\chapter{光量子、热辐射与半经典量子模型}
\section{光电效应与表面势垒}\label{th:m4photo}
电子经电压加速，撞击靶材可以把动能转化为X射线。若单个电子最多把全部动能交给一个光子，截止波长满足$hc/\lambda_{\min}=eU$。光电效应研究相反的能量转移：入射光使材料中的电子逸出。单光子模型中，光子能量为$h\nu$，电子逸出需克服功函数$W$，最大动能
\[
 K_{\max}=h\nu-W,\qquad eU_s=K_{\max}.
\]
其中$e>0$，$U_s$为使最快光电子也到不了收集极的反向电压大小。只有$\nu\geq W/h$才可逸出。提高入射功率通常增加单位时间可吸收的光子数，不会把每个光子的能量改为更大；截止频率和遏止电压随频率的变化，因而由单光子能量决定。

真实金属还存在能带占据、散射、表面状态及逸出方向分布。下面把这些细节简化为明确的势垒模型，用以讨论方向筛选；其逃逸比例不应直接当作任意材料的实测量子效率。

\begin{example}\label{ex:m4escape}
平面表面法线沿$z$正向，材料位于$z<0$。每个被吸收的光子使一个初始动能可忽略的电子获得总动能$h\nu$，速度方向在整个$4\pi$内各向同性。表面势垒只改变法向动量，势垒高度为$W$。求逃逸概率，并求收集极施加反向电压$U\geq0$时的电流。吸收光功率为$P$，忽略空间电荷和二次散射。
\Diagram{fig:m4escapeq}{电子方向角相对外法线量取；只有法向能量用于越过势垒}{
\draw[curve](-2,0)--(2,0);\node[left]at(-2,-.5){材料};
\draw[axis](0,-.7)--(0,2.5)node[above]{$z$};
\draw[->,thick](0,0)--(1.3,1.8)node[right]{$\vct v$};
\draw[guide](0,0)--(-1.3,1.8);
\draw(0,.72)arc(90:54.16:.72);\node at(.31,.92){$\theta$};
\node[right]at(2,.2){$W$};
}
\end{example}
\begin{solution}
图\ref{fig:m4escapeq}的电子法向初动能为$h\nu\cos^2\theta$。平行表面的动量守恒，不能用于克服法向势垒；逸出还必须$\cos\theta>0$。若$h\nu<W$，概率为零；否则允许锥满足$\cos\theta\geq\sqrt{W/(h\nu)}$，所以
\[
 \mathcal P_0=\frac1{4\pi}\int_0^{2\pi}\dd\varphi
 \int_0^{\theta_c}\sin\theta\dd\theta
 =\frac12[-\cos\theta]_0^{\theta_c}
 =\frac12\left(1-\sqrt{\frac W{h\nu}}\right).
\]
\Diagram{fig:m4escapea}{反向电压缩小可被收集的方向锥}{\draw[axis](0,0)--(0,2.8)node[above]{法线};\draw[curve](0,0)--(1.5,2);\draw[curve](0,0)--(-1.5,2);\draw[very thick](0,0)--(.7,2);\draw[very thick](0,0)--(-.7,2);\node[right]at(1.5,2){$U=0$};\node[right]at(.7,2.6){$U>0$};}
反向电场进一步消耗法向能量$eU$。图\ref{fig:m4escapea}中的窄锥满足$h\nu\cos^2\theta\geq W+eU$，由相同积分得
\[
 I(U)=\frac{eP}{2h\nu}
 \left[1-\sqrt{\frac{W+eU}{h\nu}}\right],\quad
 0\leq eU\leq h\nu-W.
\]
超过该区间时$I=0$。即使$h\nu\gg W$，此模型零电压的逃逸比例也只趋向$1/2$，因为另半球的电子朝材料内部。阈值电压仍为$U_s=(h\nu-W)/e$，与简单能量关系一致。
\end{solution}

\begin{example}\label{ex:m4stopping}
同一材料先后受频率$\nu_1,\nu_2$照射，测得遏止电压$U_1,U_2$，其中$\nu_2>\nu_1$。推导由测量求$h/e$和功函数的方法，并解释光电流加倍是否改变结果。
\Diagram{fig:m4stoppingq}{光电管的反向电压测量}{\draw[curve](0,0)rectangle(4,2);\draw[very thick](.8,.4)--(.8,1.6);\draw[very thick](3.2,.4)--(3.2,1.6);\draw[->](-1,1.4)--(.7,1.4)node[midway,above]{$\nu$};\node[below]at(.8,.3){阴极};\node[below]at(3.2,.3){收集极};\node[below]at(2,-.4){反向电压$U$};}
\end{example}
\begin{solution}
图\ref{fig:m4stoppingq}对每种频率分别满足$eU_i=h\nu_i-W$。两式相减，得到$h/e=(U_2-U_1)/(\nu_2-\nu_1)$，再代回$W/e=(h/e)\nu_1-U_1$。
\Diagram{fig:m4stoppinga}{遏止电压对频率的直线斜率与截距}{\draw[axis](0,0)--(4.6,0)node[right]{$\nu$};\draw[axis](0,-.7)--(0,3)node[above]{$U_s$};\draw[curve](.8,0)--(4,2.6);\fill(1.7,.731)circle(1.5pt);\fill(3.2,1.95)circle(1.5pt);\node[below]at(.8,0){$W/h$};\node[right]at(3.3,2.6){斜率$h/e$};}
图\ref{fig:m4stoppinga}的横截距给截止频率。在线性单光子区且材料表面未改变时，增大光功率只增加电子数，直线斜率和截距不变。若电流很大产生空间电荷，或进入多光子吸收区，就超出了此模型。
\end{solution}

\section{腔辐射的模态与能量分布}\label{th:m4modes}
黑体完全吸收照到它的辐射，热平衡时其辐射谱只由温度决定。小孔腔体可近似黑体：外光多次反射后被吸收，孔口射出的光取样于腔内平衡辐射。记$u_\nu(T)\dd\nu$为单位体积在频率区间内的能量，$u=\int u_\nu\dd\nu$。

\subsection{从腔内能量到孔口通量}
各向同性辐射的单位立体角能量密度为$u/(4\pi)$。在$\dd t$内，从方向$\theta$穿过面积$A$的光来自体积$Ac\dd t\cos\theta\dd\Omega$，只积分向外半球。因此出射功率面密度
\[
 M=\frac{cu}{4\pi}\int_0^{2\pi}\dd\varphi\int_0^{\pi/2}\cos\theta\sin\theta\dd\theta=\frac{cu}4.
\]
不能写成$cu/2$，因为斜向光的法向投影使平均值再多一因子$1/2$。

\subsection{三维电磁驻波计数}
取体积$V=L^3$的立方腔，先使用周期边界，波矢各分量间隔$2\pi/L$。每个波矢有两个横向极化态，半径$k$的球壳中模态数
\[
 \dd N=2\frac{V}{(2\pi)^3}4\pi k^2\dd k.
\]
真空$k=2\pi\nu/c$，所以单位体积、单位频率的态密度
\begin{equation}\label{eq:m4density}
 g(\nu)=\frac1V\frac{\dd N}{\dd\nu}=\frac{8\pi\nu^2}{c^3}.
\end{equation}
若改用理想导体驻波边界，正整数波数间隔为$\pi/L$，只取球壳八分之一，再乘两极化，宏观极限结果相同。边界附近少数特殊模态只影响低阶有限尺寸修正。计数时不可把周期边界的全空间与驻波边界的八分球重复叠加。

\Diagram{fig:m4modes}{周期边界波矢格点与球壳；每点有两种横向极化}{\draw[axis](0,0)--(4,0)node[right]{$k_x$};\draw[axis](0,0)--(0,3.6)node[above]{$k_y$};\foreach\x in{.5,1,1.5,2,2.5,3}{\foreach\y in{.5,1,1.5,2,2.5,3}{\fill(\x,\y)circle(.7pt);}}\draw[curve](3,0)arc(0:90:3);\draw[curve](3.3,0)arc(0:90:3.3);\node[right]at(2.8,2){$k\to k+\dd k$};}

\section{从经典配分到普朗克谱}\label{th:m4planck}
经典谐振模态有两个二次型能量项，平均各为$\kb T/2$，所以每模平均总能为$\kb T$。也可把模态的相空间按能量积分：谐振子的等能椭圆所围面积正比能量，能量态密度为常数，故
\[
\langle E\rangle=\frac{\int_0^\infty E\ee^{-E/(\kb T)}\dd E}{\int_0^\infty\ee^{-E/(\kb T)}\dd E}.
\]
令$b=\kb T$。分母为$[-b\ee^{-E/b}]_0^\infty=b$；分子分部积分为$[-bE\ee^{-E/b}]_0^\infty+b\int_0^\infty\ee^{-E/b}\dd E=b^2$，故$\langle E\rangle=b=\kb T$。
若只积分一个动量变量的高斯分布，得到的是一个二次项的$\kb T/2$，不能把它当成整模能量。经典谱$u_\nu=8\pi\nu^2\kb T/c^3$在高频积分发散，这称为紫外灾难。

取单模可交换能量为$E_n=nh\nu$，$n=0,1,\ldots$。令$q=\ee^{-h\nu/(\kb T)}$，归一化配分和为$Z=\sum q^n=1/(1-q)$。利用$q\dd Z/\dd q=\sum nq^n=q/(1-q)^2$，得
\[
 \langle E\rangle=h\nu\frac q{1-q}=\frac{h\nu}{\ee^{h\nu/(\kb T)}-1}.
\]
乘态密度得到普朗克频率谱
\begin{equation}\label{eq:m4planck}
 u_\nu=\frac{8\pi h\nu^3}{c^3}\frac1{\ee^{h\nu/(\kb T)}-1}.
\end{equation}
低频$h\nu\ll\kb T$时展开指数，恢复经典谱；高频端指数压低，积分收敛。谐振子总能量还可含$h\nu/2$零点项，但这一与温度无关的基态能不作为黑体热辐射可交换能量计入此处。

\begin{example}\label{ex:m4stefan}
由普朗克谱推导$M=\sigma T^4$，并给出$\sigma$。积分中可使用$\sum_{n=1}^\infty n^{-4}=\pi^4/90$，但需说明它怎样进入结果。
\Diagram{fig:m4stefanq}{腔内平衡辐射经小孔向外半球射出}{\draw[curve](0,0)--(0,3)--(3,3)--(3,1.8);\draw[curve](3,1.2)--(3,0)--(0,0);\draw[->](3,1.5)--(5,2.4);\draw[->](3,1.5)--(5,1.5);\draw[->](3,1.5)--(5,.6);\node at(1.5,1.5){$u(T)$};\node[right]at(5,1.5){$M$};}
\end{example}
\begin{solution}
令$x=h\nu/(\kb T)$，由式\eqref{eq:m4planck}，
\[
 u=\frac{8\pi\kb^4T^4}{h^3c^3}\int_0^\infty\frac{x^3\dd x}{\ee^x-1}.
\]
对$x>0$，$(\ee^x-1)^{-1}=\sum_{n=1}^\infty\ee^{-nx}$。各项非负，允许交换求和积分。令$y=nx$，则$\dd x=\dd y/n$、上下限仍为$0,\infty$，每一项为$n^{-4}\int_0^\infty y^3\ee^{-y}\dd y$。分部积分取$\dd v=\ee^{-y}\dd y$、$v=-\ee^{-y}$，依次得
\begin{align*}
 \int_0^\infty y^3\ee^{-y}\dd y
 &=[-y^3\ee^{-y}]_0^\infty+3\int_0^\infty y^2\ee^{-y}\dd y,\\
 \int_0^\infty y^2\ee^{-y}\dd y
 &=[-y^2\ee^{-y}]_0^\infty+2\int_0^\infty y\ee^{-y}\dd y,\\
 \int_0^\infty y\ee^{-y}\dd y
 &=[-y\ee^{-y}]_0^\infty+\int_0^\infty\ee^{-y}\dd y=1.
\end{align*}
三个边界项都为零，最后一个积分为$[-\ee^{-y}]_0^\infty=1$，所以三个积分依次为$6,2,1$。因此
\[
 \int_0^\infty\frac{x^3\dd x}{\ee^x-1}=6\sum_{n=1}^\infty n^{-4}=\frac{\pi^4}{15}.
\]
\Diagram{fig:m4stefana}{角度投影积分把腔内$cu$变为孔口$cu/4$}{\draw[curve](-1.8,0)--(1.8,0);\draw[axis](0,0)--(0,2.8)node[above]{法线};\draw[axis](0,0)--(1.5,2)node[right]{$c\dd t$};\draw[guide](1.5,0)--(1.5,2)--(0,2);\node[left]at(0,1){$c\dd t\cos\theta$};}
图\ref{fig:m4stefanq}的孔口通量按图\ref{fig:m4stefana}作半球投影积分：
\[
 M=\frac{cu}{4\pi}\int_0^{2\pi}\dd\varphi
 \int_0^{\pi/2}\cos\theta\sin\theta\dd\theta
 =\frac{cu}{2}\left[\frac{\sin^2\theta}{2}\right]_0^{\pi/2}
 =\frac{cu}{4}.
\]
代入能量密度，得到
\[
 \sigma=\frac{2\pi^5\kb^4}{15h^3c^2}.
\]
因$u$是体积能量密度，$cu$的量纲为功率每面积，与$M$一致。温度的四次方来自三次频率因子和频率积分的一次因子。
\end{solution}

\begin{example}\label{ex:m4wien}
分别求频率谱$u_\nu$和波长谱$u_\lambda$峰值对应的无量纲方程，说明为何$\lambda_{\max}\nu_{\max}\ne c$。
\Diagram{fig:m4wienq}{相同能量小区间在频率轴和波长轴上宽度不同}{\draw[axis](0,1.5)--(5,1.5)node[right]{$\nu$};\draw[axis](0,0)--(5,0)node[right]{$\lambda$};\draw[very thick](1,1.5)--(2,1.5);\draw[very thick](2,0)--(4,0);\draw[guide](1,1.5)--(4,0);\draw[guide](2,1.5)--(2,0);}
\end{example}
\begin{solution}
图\ref{fig:m4wienq}中的频率区间与波长区间方向相反。能量相等应写$u_\lambda\dd\lambda=-u_\nu\dd\nu$，因$\nu=c/\lambda$，
\[
 u_\lambda=u_\nu(c/\lambda)\left|\frac{\dd\nu}{\dd\lambda}\right|
 =\frac{8\pi hc}{\lambda^5[\ee^{hc/(\lambda\kb T)}-1]}.
\]
频率谱中令$x=h\nu/(\kb T)$，只需极大化$f_3(x)=x^3/(\ee^x-1)$。对一般$f_j(q)=q^j/(\ee^q-1)$使用商法则，
\[
 f_j'(q)=\frac{jq^{j-1}(\ee^q-1)-q^j\ee^q}{(\ee^q-1)^2}
 =\frac{q^{j-1}\ee^q}{(\ee^q-1)^2}
 [j(1-\ee^{-q})-q].
\]
令$j=3$，非零峰便满足$3(1-\ee^{-x})=x$。
\Diagram{fig:m4wiena}{频率谱和波长谱的峰值各由自己的密度函数决定}{\draw[axis](0,0)--(5,0)node[right]{无量纲变量};\draw[axis](0,0)--(0,2.6)node[above]{归一化谱};\draw[curve,domain=.1:7,samples=90]plot({.6*\x},{1.4*\x^3/(exp(\x)-1)});\draw[guide](1.693,0)--(1.693,2);\node[below]at(1.693,0){$2.8214$};\node[right]at(4.3,1.8){$u_\nu$示意};}
图\ref{fig:m4wiena}的频率峰根为$x_\nu\simeq2.82144$。对波长谱令$y=hc/(\lambda\kb T)$，它正比$f_5(y)$。链式法则给$\dd u_\lambda/\dd\lambda\propto f_5'(y)\dd y/\dd\lambda$，其中$\dd y/\dd\lambda=-hc/(\kb T\lambda^2)\ne0$。在上式取$j=5$，峰条件为$5(1-\ee^{-y})=y$，非零正根$y_\lambda\simeq4.96511$。因此
\[
 \nu_{\max}=\frac{x_\nu\kb T}h,\quad
 \lambda_{\max}=\frac{hc}{y_\lambda\kb T},\quad
 \lambda_{\max}\nu_{\max}=c\frac{x_\nu}{y_\lambda}\simeq.56825c.
\]
两个方程都有形式根0，但谱在原点或相应无穷远边界趋零，它不是所求内部峰。两谱各自只有一个非零正峰；其位置不同来自密度换元的雅可比$c/\lambda^2$，不表示同一个光子不满足$\lambda\nu=c$。
\end{solution}

\section{玻尔氢样原子与相对论圆轨道}\label{th:m4bohr}
固定核电荷$+Ze$，电子质量$m$，忽略核反冲。非相对论圆轨道满足$m v^2/r=Ze^2/(4\pi\varepsilon_0r^2)$、$mvr=n\hbar$。消去$v$得到
\[
 r_n=\frac{a_0n^2}Z,\quad v_n=\frac{Z\alpha c}n,\quad
 E_n=-\frac{mc^2(Z\alpha)^2}{2n^2},
\]
其中$\alpha=e^2/(4\pi\varepsilon_0\hbar c)$、$a_0=\hbar/(mc\alpha)$，$n=1,2,\ldots$。总能量等于动能加库仑势能，利用向心关系势能为动能的$-2$倍，所以能量为负。长度$a_0$、动量$m\alpha c$、能量$m\alpha^2c^2$组成一组方便的原子单位；氢基态束缚能为这一能量单位的一半。

光谱跃迁频率由$h\nu=E_{n_i}-E_{n_f}>0$确定。圆轨道的转动频率只是大量子数极限下相邻能级间隔的近似，不可对任意$n_i,n_f$直接取轨道频率作为谱线频率。

\begin{example}\label{ex:m4relbohr}
保留圆轨道角动量量子化$pr=n\hbar$，把动力学改为相对论，库仑势保持$-Ze^2/(4\pi\varepsilon_0r)$。求$v_n,r_n$、包含静能的总能量及其低$Z\alpha/n$展开。
\Diagram{fig:m4relbohrq}{圆轨道中动量切向、库仑力径向}{\draw[curve](0,0)circle(1.5);\fill(0,0)circle(2pt)node[below]{核};\fill(1.5,0)circle(2pt);\draw[axis](1.5,0)--(1.5,1.3)node[right]{$p$};\draw[axis](1.5,0)--(.3,0)node[above]{$F$};\node[below]at(.75,0){$r$};}
\end{example}
\begin{solution}
图\ref{fig:m4relbohrq}中动量大小不变但方向转动，径向力为$pv/r$，故$pv=Ze^2/(4\pi\varepsilon_0r)$。代入$pr=n\hbar$，得到$v/c=Z\alpha/n\equiv\eta$。于是$p=\gamma_vmv$，半径为
\[
 r=\frac{n\hbar}{\gamma_vmv}=\frac{a_0n^2}Z\sqrt{1-\eta^2}.
\]
包含静能的总能量$\mathcal E=\gamma_vmc^2-Ze^2/(4\pi\varepsilon_0r)$。利用势能绝对值$pv=\gamma_vmv^2$，得
\[
 \mathcal E=\gamma_vmc^2(1-\eta^2)=mc^2\sqrt{1-\eta^2}.
\]
\Diagram{fig:m4relbohra}{相对论圆轨道的束缚能随$\eta$增加而加深}{\draw[axis](0,0)--(4.2,0)node[right]{$\eta$};\draw[axis](0,-2.6)--(0,.7)node[above]{$E_b/(mc^2)$};\draw[curve,domain=0:.98,samples=70]plot({3.6*\x},{2.3*(sqrt(1-\x*\x)-1)});\draw[guide,domain=0:1,samples=40]plot({3.6*\x},{-1.15*\x*\x});}
令$z=\eta^2$，对$f(z)=\sqrt{1-z}$求导有$f'(0)=-1/2$、$f''(0)=-1/4$。泰勒展开$f(z)=1+f'(0)z+f''(0)z^2/2+O(z^3)=1-z/2-z^2/8+O(z^3)$。减去静能后，图\ref{fig:m4relbohra}的束缚能为
\[
 E_b=-mc^2\left(\frac{\eta^2}2+\frac{\eta^4}8+O(\eta^6)\right).
\]
首项恢复玻尔能级。模型要求$\eta<1$即$n>Z\alpha$。它只修正半经典圆轨道，不含自旋、非圆轨道的角动量结构与完整量子动力学，因此不能把此式当作全部精细结构能级。
\end{solution}

\section{作用量量子化与转折边界}\label{th:m4action}
对一维束缚运动，定义一个完整周期的作用量$J=\oint p\dd x$。这是相平面闭合轨道所围面积。圆周运动则$J_\varphi=\int_0^{2\pi}L\dd\varphi=2\pi L$，旧量子条件$J_\varphi=nh$给$L=n\hbar$。

一维波函数在允许区的短波形式含相位$\int p\dd x/\hbar$。边界反射时还会附加相位：两端硬壁使$J=nh$，$n=1,2,\ldots$；两端平滑转折点使$J=(n+1/2)h$，$n=0,1,\ldots$；一硬壁、一平滑转折点使$J=(n+3/4)h$。后两式是半经典近似，一般在量子数较大时更可靠。谐振子恰能由该条件给出精确能级，但不能据此认定任意势的半经典最低能级也精确。

\begin{example}\label{ex:m4well}
质量$m$的粒子限制在$0<x<L$，外侧无限高势垒，内部势能为零。由驻波与作用量两种方法求能级，并写出归一化波函数。
\Diagram{fig:m4wellq}{无限深方阱的两个硬边界}{\draw[axis](-.5,0)--(4.8,0)node[right]{$x$};\draw[curve](0,3)--(0,0)--(3.5,0)--(3.5,3);\node[below]at(0,0){$0$};\node[below]at(3.5,0){$L$};\node[above]at(1.75,2.1){$V=0$};}
\end{example}
\begin{solution}
图\ref{fig:m4wellq}的阱内满足$-\hbar^2\psi''/(2m)=E\psi$。令$k=\sqrt{2mE}/\hbar$，通解为$A\sin kx+B\cos kx$。左壁$\psi(0)=0$给$B=0$；右壁$\psi(L)=0$要求$kL=n\pi$，$n=1,2,\ldots$。$n=0$只给处处为零的函数，不能归一化，故不属于允许态。
\[
 E_n=\frac{\hbar^2k_n^2}{2m}=\frac{n^2h^2}{8mL^2}.
\]
\Diagram{fig:m4wella}{方阱相轨道为两条水平动量支，面积为$2pL$}{\draw[axis](-.3,0)--(4.6,0)node[right]{$x$};\draw[axis](0,-1.8)--(0,2)node[above]{$p$};\draw[very thick,->](0,1.2)--(3.5,1.2);\draw[very thick,->](3.5,-1.2)--(0,-1.2);\draw[guide](0,1.2)--(0,-1.2);\draw[guide](3.5,1.2)--(3.5,-1.2);\node[right]at(3.5,1.2){$p_n$};}
归一化积分逐步为
\[
 1=|A|^2\int_0^L\frac{1-\cos(2n\pi x/L)}2\dd x
 =|A|^2\left[\frac x2-\frac{L}{4n\pi}\sin\frac{2n\pi x}{L}\right]_0^L
 =|A|^2L/2,
\]
故可取$\psi_n=\sqrt{2/L}\sin(n\pi x/L)$。图\ref{fig:m4wella}的相轨道上，向右支$p=p_n$、$\dd x>0$；向左支$p=-p_n$、$\dd x<0$，因此两支贡献都为正：
\[
 J=\int_0^L p_n\dd x+\int_L^0(-p_n)\dd x
 =[p_nx]_0^L+[-p_nx]_L^0=2p_nL.
\]
用两个硬壁条件$J=nh$同样得到$p_n=nh/(2L)$。竖直碰壁段满足$\dd x=0$，不增加作用量。这个定态是两个反向行波的叠加，$\langle p\rangle=0$，却有$\langle p^2\rangle=p_n^2$；不能将相轨道上的固定动量大小解释为定态粒子总沿某一个方向运动。
\end{solution}

\begin{example}\label{ex:m4osc}
一维谐振子势$V(x)=m\omega^2x^2/2$。用相平面面积求半经典能级，并比较直接使用$J=nh$会遗漏什么。
\Diagram{fig:m4oscq}{谐振子的两个平滑转折点}{\draw[axis](-2.3,0)--(2.5,0)node[right]{$x$};\draw[axis](0,0)--(0,3)node[above]{$V$};\draw[curve,domain=-1.9:1.9,samples=60]plot(\x,{.65*\x*\x});\draw[guide](-1.7,1.88)--(1.7,1.88);\node[right]at(1.7,1.88){$E$};}
\end{example}
\begin{solution}
图\ref{fig:m4oscq}的势为$V=m\omega^2x^2/2$。固定能量的相轨道满足
\[
 \frac{x^2}{2E/(m\omega^2)}+\frac{p^2}{2mE}=1,
\]
故两半轴$x_m=\sqrt{2E/(m\omega^2)}$、$p_m=\sqrt{2mE}$。作用量是顺运动方向的闭合积分，上下两支都给正面积贡献。
\Diagram{fig:m4osca}{椭圆所围面积$J=\pi x_mp_m$}{\draw[axis](-2.6,0)--(2.8,0)node[right]{$x$};\draw[axis](0,-1.9)--(0,2.2)node[above]{$p$};\draw[curve](0,0)ellipse(2 and 1.4);\node[below right]at(2,0){$x_m$};\node[above right]at(0,1.4){$p_m$};\draw[->](1.8,.61)--(1.95,.31);}
图\ref{fig:m4osca}的椭圆面积直接给$J=\pi x_mp_m=2\pi E/\omega$。也可令$x=x_m\sin\varphi$作积分：
\[
 J=4p_mx_m\int_0^{\pi/2}\cos^2\varphi\dd\varphi
 =4p_mx_m\left[\frac\varphi2+\frac{\sin2\varphi}{4}\right]_0^{\pi/2}
 =\pi p_mx_m.
\]
两个转点平滑，取$J=(n+1/2)h$，$n=0,1,\ldots$，于是$E_n=\hbar\omega(n+1/2)$，相邻能级间隔为$\hbar\omega$。对谐振子这一量子条件恰给出精确能级，但不能据此认定任意势的最低能级都由半经典条件精确给出。

若用$J=nh$并允许$n=0$，便给出$E=0$、$x=p=0$，与量子态不能同时具有确定位置和动量的要求不符。零点项来自转点相位，并非任意选取能量零点；给定势最低点为零后，基态相对最低点仍有$\hbar\omega/2$。由$\dd J/\dd E=2\pi/\omega$还恢复能量无关的经典周期。
\end{solution}

\begin{example}\label{ex:m4bounce}
粒子在$x\geq0$受重力势$V=mgx$，$x<0$为无限高硬壁。用一硬一软的量子条件求能级近似，并说明其量纲与随量子数的变化。
\Diagram{fig:m4bounceq}{底部硬壁与上方重力转折点}{\draw[axis](-.4,0)--(4.5,0)node[right]{$x$};\draw[axis](0,0)--(0,3.2)node[above]{$V$};\draw[very thick](0,0)--(0,3);\draw[curve](0,0)--(3.8,2.8);\draw[guide](0,2)--(2.714,2);\node[right]at(2.8,2){$x_t=E/(mg)$};}
\end{example}
\begin{solution}
图\ref{fig:m4bounceq}中的允许区是$0\leq x\leq E/(mg)$，动量大小$p(x)=\sqrt{2m(E-mgx)}$。完整往返作用量
\begin{align*}
 J&=2\int_0^{E/(mg)}\sqrt{2m(E-mgx)}\dd x\\
 &=-\frac{2\sqrt{2m}}{mg}\int_E^0 q^{1/2}\dd q
 =\frac{2\sqrt{2m}}{mg}\left[\frac23q^{3/2}\right]_0^E
 =\frac{4\sqrt2}{3g\sqrt m}E^{3/2}.
\end{align*}
这里取$q=E-mgx$，故$\dd x=-\dd q/(mg)$，空间上下限$0,E/(mg)$依次变为$E,0$。
\Diagram{fig:m4bouncea}{上下动量支在平滑转点相连，在硬壁处反射}{\draw[axis](-.3,0)--(4.5,0)node[right]{$x$};\draw[axis](0,-2)--(0,2.2)node[above]{$p$};\draw[curve,domain=0:3.5,samples=70]plot(\x,{1.6*sqrt(1-\x/3.5)});\draw[curve,domain=0:3.5,samples=70]plot(\x,{-1.6*sqrt(1-\x/3.5)});\draw[very thick](0,-1.6)--(0,1.6);\node[below]at(3.5,0){$x_t$};}
图\ref{fig:m4bouncea}中一端是硬壁，一端是平滑转点，故取$J=(n+3/4)h$，得到
\[
 E_n=\left[\frac{3hg\sqrt m}{4\sqrt2}\left(n+\frac34\right)\right]^{2/3},\quad n=0,1,\ldots.
\]
括号内量纲为能量的$3/2$次方。能量随$n^{2/3}$增长，所以高能级的相邻间隔减小。此式是半经典近似；精确解应以满足硬壁条件的艾里函数零点确定，不能把$3/4$条件称为最低几个能级的精确值。
\end{solution}
\section{从波长关系到束缚态方程}
\subsection{经典轨迹和量子态的不同描述}
玻尔模型把经典圆轨道与量子化条件相结合，能解释氢样原子的主要能量尺度，却没有给出一般波函数和测量概率。处理波函数节点、势垒穿透和一般角动量态，需要量子态的描述。更完整的一维描述使用复波函数$\psi(x,t)$，$|\psi|^2$为位置概率密度，满足归一化$\int|\psi|^2\dd x=1$。

自由粒子的德布罗意关系$p=\hbar k$、$E=\hbar\omega$与非相对论色散$E=p^2/(2m)$相容于方程$i\hbar\partial_t\psi=-\hbar^2\partial_x^2\psi/(2m)$。加入给定势能$V(x)$，非相对论量子理论采用
\[
 i\hbar\partial_t\psi=\left[-\frac{\hbar^2}{2m}\partial_x^2+V(x)\right]\psi.
\]
这是一条量子动力学的基本方程；前面的自由粒子关系只说明它与德布罗意关系相容。对不含时势，取$\psi=\varphi(x)\ee^{-iEt/\hbar}$，则$\partial_t\psi=(-iE/\hbar)\psi$、$\partial_x^2\psi=\varphi''\ee^{-iEt/\hbar}$。代入并约去共同的时间因子，得到$-\hbar^2\varphi''/(2m)+V\varphi=E\varphi$。允许能量还受边界条件约束，不能只由局部$p^2=2m(E-V)$确定。

\subsection{边界连续性与允许区}
有限势跳变处，积分定态方程跨过一个无穷小区间。只要势没有狄拉克型奇点，积分中的有限项趋零，得到导数连续；波函数本身也连续，避免二阶导数产生未被势抵消的奇异项。无限高硬壁是极限边界，波函数在壁外为零、壁处为零，壁内导数不需要与外侧零函数继续相等。

在$E>V$区，局部波数实，波函数可振荡；$E<V$区，局部波数纯虚，波函数可以指数衰减或增长。有限势垒内的非零尾部不表示经典粒子拥有负动能，而是波函数满足另一类局部微分方程。对半无限禁区排除增长解，对有限禁区则一般保留两支以匹配两端。

\Diagram{fig:m4allowed}{线性势的艾里解平滑连接允许区与禁区}{\draw[axis](-.4,0)--(6,0)node[right]{$x$};\draw[axis](0,-1.2)--(0,1.9)node[above]{$\varphi$};\draw[guide](3,-1.2)--(3,1.8);\draw[curve]plot coordinates {(0.00000,-0.9470357) (0.05046,-0.9038078) (0.10092,-0.8537957) (0.15138,-0.7974842) (0.20183,-0.7353909) (0.25229,-0.6680602) (0.30275,-0.5960570) (0.35321,-0.5199612) (0.40367,-0.4403617) (0.45413,-0.3578511) (0.50459,-0.2730203) (0.55505,-0.1864540) (0.60550,-0.0987258) (0.65596,-0.0103941) (0.70642,0.0780016) (0.75688,0.1659438) (0.80734,0.2529406) (0.85780,0.3385279) (0.90826,0.4222713) (0.95872,0.5037687) (1.00917,0.5826506) (1.05963,0.6585822) (1.11009,0.7312630) (1.16055,0.8004281) (1.21101,0.8658473) (1.26147,0.9273255) (1.31193,0.9847018) (1.36239,1.0378485) (1.41284,1.0866705) (1.46330,1.1311037) (1.51376,1.1711138) (1.56422,1.2066944) (1.61468,1.2378660) (1.66514,1.2646733) (1.71560,1.2871841) (1.76606,1.3054872) (1.81651,1.3196902) (1.86697,1.3299179) (1.91743,1.3363102) (1.96789,1.3390201) (2.01835,1.3382122) (2.06881,1.3340601) (2.11927,1.3267457) (2.16972,1.3164565) (2.22018,1.3033846) (2.27064,1.2877250) (2.32110,1.2696741) (2.37156,1.2494285) (2.42202,1.2271835) (2.47248,1.2031323) (2.52294,1.1774648) (2.57339,1.1503667) (2.62385,1.1220186) (2.67431,1.0925954) (2.72477,1.0622656) (2.77523,1.0311908) (2.82569,0.9995253) (2.87615,0.9674155) (2.92661,0.9349999) (2.97706,0.9024089) (3.02752,0.8697645) (3.07798,0.8371804) (3.12844,0.8047618) (3.17890,0.7726057) (3.22936,0.7408008) (3.27982,0.7094279) (3.33028,0.6785597) (3.38073,0.6482616) (3.43119,0.6185912) (3.48165,0.5895995) (3.53211,0.5613302) (3.58257,0.5338209) (3.63303,0.5071029) (3.68349,0.4812018) (3.73394,0.4561376) (3.78440,0.4319255) (3.83486,0.4085756) (3.88532,0.3860938) (3.93578,0.3644820) (3.98624,0.3437383) (4.03670,0.3238576) (4.08716,0.3048314) (4.13761,0.2866488) (4.18807,0.2692963) (4.23853,0.2527582) (4.28899,0.2370171) (4.33945,0.2220536) (4.38991,0.2078474) (4.44037,0.1943766) (4.49083,0.1816186) (4.54128,0.1695500) (4.59174,0.1581466) (4.64220,0.1473842) (4.69266,0.1372379) (4.74312,0.1276832) (4.79358,0.1186951) (4.84404,0.1102491) (4.89450,0.1023207) (4.94495,0.0948859) (4.99541,0.0879210) (5.04587,0.0814028) (5.09633,0.0753088) (5.14679,0.0696167) (5.19725,0.0643052) (5.24771,0.0593535) (5.29817,0.0547415) (5.34862,0.0504499) (5.39908,0.0464601) (5.44954,0.0427541) (5.50000,0.0393148)};\node[above]at(1.2,1.8){允许区};\node[above]at(4.5,1.8){禁区};}
图\ref{fig:m4allowed}画出线性势中的同一个艾里函数解：转点左侧振荡，右侧衰减，函数及其导数在转点处连续。两区振幅由这个连接条件确定，不能分别任取。有限方阱解中的$\arctan(k/\kappa)$边界相位同样来自函数和导数的同时匹配。

\section{WKB相位、振幅和量子偏移}
令定态波函数$\varphi=A\ee^{iS/\hbar}$。先求一次导数$\varphi'=(A'+iAS'/\hbar)\ee^{iS/\hbar}$，再对括号与指数分别求导，得到
\[
 \varphi''=\left[A''+\frac{2i}{\hbar}A'S'
 +\frac{i}{\hbar}AS''-\frac1{\hbar^2}A(S')^2\right]\ee^{iS/\hbar}.
\]
当$A,S$相对局部波长缓变，领先阶给$(S')^2=2m(E-V)=p^2(x)$；下一阶虚部给$2A'p+Ap'=0$。在$p>0$的一个允许区间内，以$x_0$为参考点分离变量，
\[
 \frac{A'}A=-\frac{p'}{2p},\qquad
 \int_{A(x_0)}^{A(x)}\frac{\dd a}{a}
 =-\frac12\int_{p(x_0)}^{p(x)}\frac{\dd q}{q},\qquad
 \ln\frac{A(x)}{A(x_0)}=-\frac12\ln\frac{p(x)}{p(x_0)}.
\]
因此$A=C/\sqrt p$，允许区的两支为
\[
 \varphi_\pm\simeq\frac{C_\pm}{\sqrt{p(x)}}
 \exp\left[\pm\frac i\hbar\int^x p(q)\dd q\right].
\]
检验忽略$A''$的条件，还须把振幅代回原方程。对$A=Cp^{-1/2}$连续求导，
\[
 A'=-\frac C2p^{-3/2}p',\qquad
 A''=\frac{3C}{4}p^{-5/2}(p')^2-\frac C2p^{-3/2}p'',
\]
故相对于领先项$p^2A$的余项为
\[
 \frac{\hbar^2A''}{p^2A}
 =\hbar^2\left[\frac{3(p')^2}{4p^4}-\frac{p''}{2p^3}\right].
\]
一组充分的局部缓变条件是$\hbar|p'|/p^2\ll1$及$\hbar^2|p''|/p^3\ll1$。长区间的累计误差还需结合具体势估计。在转点$p=0$附近条件失效，不能把上述振幅直接延伸穿过转点。

\subsection{平滑转点的四分之一周期相移}
取右侧简单转点$x_t$，即$V(x_t)=E$、$F_t=V'(x_t)>0$，允许区在$x<x_t$。局部线性化$V-E\simeq F_t(x-x_t)$，定态方程成为$\varphi''=(2mF_t/\hbar^2)(x-x_t)\varphi$。令$\ell_t=(\hbar^2/(2mF_t))^{1/3}$、$z=(x-x_t)/\ell_t$，则$\varphi''(x)=\ell_t^{-2}\varphi_{zz}$，代入得$\varphi_{zz}-z\varphi=0$。禁区$z>0$选衰减的$\operatorname{Ai}(z)$。利用艾里函数在$z<0$的渐近连接式，允许区的驻波为
\[
 \varphi(x)\propto\frac1{\sqrt{p(x)}}
 \sin\left[\frac1\hbar\int_x^{x_t}p(q)\dd q+\frac\pi4\right].
\]
这里的相移是局部线性势的连接结果，不能用于突然跳到无限高的硬壁。两个平滑转点各带相同类型修正，闭合相位条件成为$\int_{x_1}^{x_2}p\dd x=(n+1/2)\pi\hbar$，即$J=(n+1/2)h$。

若左端是硬壁$x=0$、右端是平滑转点，将上式连接解代入$\varphi(0)=0$，得到
\[
 \frac1\hbar\int_0^{x_t}p\dd x+\frac\pi4=(n+1)\pi,
\]
所以$J=2\int_0^{x_t}p\dd x=(n+3/4)h$。这给出重力反弹中的$3/4$。两硬壁之间则直接要求站波节点，$\int_0^Lp\dd x=n\pi\hbar$，$J=nh$，从$n=1$开始。三类边界的量子数起点和偏移不同，比较时须把编号一并说明。

\Diagram{fig:m4boundarytypes}{硬壁节点与平滑转点连接所对应的相位条件不同}{\begin{scope}[xshift=-2.4cm]\draw[axis](0,0)--(2.2,0);\draw[very thick](0,-.4)--(0,1.7);\draw[curve,domain=0:2,samples=60]plot(\x,{.9*sin(120*\x)});\node[below]at(1,-.7){硬壁节点};\end{scope}\begin{scope}[xshift=1.2cm]\draw[axis](0,0)--(2.5,0);\draw[curve](0,.7)..controls(.5,1.3)and(.7,1.1)..(1,.5)..controls(1.4,.1)and(1.8,.02)..(2.3,0);\draw[guide](1,-.4)--(1,1.7);\node[below]at(1.2,-.7){平滑转点};\end{scope}}
图\ref{fig:m4boundarytypes}仅区分边界形态。半经典能级公式通常对高量子数较准，因为大部分允许区包含许多波长，转点局部修正占比小。重力反弹可用艾里零点检验误差，谐振子则因二次势的特殊结构使领先WKB能级恰好精确。精确的一个例子不能消除近似条件。

\section{零点能和位置—动量展宽}
对归一化且边界项消失的波函数，定义$\Delta x^2=\langle(x-\langle x\rangle)^2\rangle$、$\Delta p^2=\langle(p-\langle p\rangle)^2\rangle$。取$f=(x-\langle x\rangle)\psi$、$g=(p-\langle p\rangle)\psi$，柯西不等式给$\Delta x^2\Delta p^2\geq|\langle f,g\rangle|^2$。由算符的厄米性及对易关系，
\[
 2i\,\operatorname{Im}\langle f,g\rangle
 =\langle f,g\rangle-\langle g,f\rangle
 =\langle\psi,[x,p]\psi\rangle=i\hbar.
\]
于是$|\langle f,g\rangle|^2\geq(\operatorname{Im}\langle f,g\rangle)^2=\hbar^2/4$，得到
\[
 \Delta x\Delta p\geq\frac\hbar2.
\]
这里的$\Delta x,\Delta p$是同一量子态的概率标准差。位置分布越窄，动量分布的最小展宽越大；方阱基态的非零动能和谐振子零点能都满足这一约束。

对以平衡点为中心、平均动量为零的谐振子态，能量平均满足
\[
 \langle E\rangle=\frac{\Delta p^2}{2m}+\frac{m\omega^2\Delta x^2}2
 \geq\frac{\hbar^2}{8m\Delta x^2}+\frac{m\omega^2\Delta x^2}2.
\]
记$z=\Delta x^2>0$，右侧函数为$f(z)=\hbar^2/(8mz)+m\omega^2z/2$。逐项求导，
\[
 f'(z)=-\frac{\hbar^2}{8mz^2}+\frac{m\omega^2}{2},\qquad
 f''(z)=\frac{\hbar^2}{4mz^3}>0.
\]
令$f'=0$，正根为$z=\hbar/(2m\omega)$，代回得$\langle E\rangle\geq\hbar\omega/2$。可取$\varphi=C\ee^{-b x^2/2}$核实等号态。两次求导为$\varphi'=-bx\varphi$、$\varphi''=(-b+b^2x^2)\varphi$，代入定态方程给
\[
 E\varphi=\left[\frac{\hbar^2b}{2m}
 +\left(\frac{m\omega^2}{2}-\frac{\hbar^2b^2}{2m}\right)x^2\right]\varphi.
\]
定态能量不依赖$x$，故$b=m\omega/\hbar$（取正根以使无穷远衰减），$E=\hbar\omega/2$。零点项同时由量子展宽下界和满足边界条件的波函数确定。

\section{黑体谱中经典近似的边界}
经典均分定理的前提是能量可在$\kb T$尺度上近似连续地取值。对低频模态，能量间隔$h\nu\ll\kb T$，相邻占据数许多，几何分布的离散和可近似为连续积分，于是每模平均$\kb T$。对高频模态，第一次激发就需远大于$\kb T$的能量，绝大多数模态停留在$n=0$，平均能量被玻尔兹曼因子压低。

紫外灾难的根源因此是把低频的连续能量近似用于全部无限多个高频模态。修正只改模态计数是不够的：真空电磁波的态密度$\nu^2$来自空间维数与色散；量子化改变的是各模态的热占据。计算时分别求出态密度和平均热占据，再将二者相乘。

黑体辐射的总能量按$T^4$、光子数按$T^3$变化，平均光子能量因而正比于$T$。把频率谱的横轴换为$h\nu/(\kb T)$、纵轴除以$T^3$，不同温度的曲线重合。实际物体若有选择性吸收、反射、有限模态或非平衡照明，其谱形可偏离这一关系。

\subsection{作用量中的反射相位与量子数编号}
一维束缚运动的作用量为$J(E)=\oint p\dd x=2\int_{x_-}^{x_+}|p(x)|\dd x$。返回支的$p$与$\dd x$同时反号，两支对作用量均给出正贡献。波的闭合相位还包括转点反射相位，因此不同边界不能统一写成$J=nh$。

两个平滑转点时，每个转点引入四分之一周的相位修正，量子条件为$J=(n+1/2)h$，$n=0,1,\ldots$。一硬一软边界给$J=(n+3/4)h$，仍从$n=0$编号；两个无限硬壁则给$J=nh$，应从$n=1$编号。三个式子的$n$从不同数值起算，分别对应各边界问题的节点编号。
\Diagram{fig:m4boundaryphase}{允许区的三类端点；硬壁与平滑转点采用不同反射条件}{\begin{scope}\draw[axis](0,0)--(1.8,0);\draw[curve]plot[smooth]coordinates{(0,1.4)(.3,.5)(.7,.1)(1.1,.5)(1.4,1.4)};\draw[guide](0,.8)--(1.5,.8);\node[below]at(.7,-.2){两软};\end{scope}\begin{scope}[xshift=2.2cm]\draw[axis](0,0)--(1.8,0);\draw[thick](.2,0)--(.2,1.5);\draw[curve](.2,.1)--(1.5,1.4);\draw[guide](.2,.8)--(1.6,.8);\node[below]at(.8,-.2){一硬一软};\end{scope}\begin{scope}[xshift=4.5cm]\draw[axis](0,0)--(1.8,0);\draw[thick](.2,1.5)--(.2,.1)--(1.4,.1)--(1.4,1.5);\draw[guide](.2,.8)--(1.4,.8);\node[below]at(.8,-.2){两硬};\end{scope}}
图\ref{fig:m4boundaryphase}的虚线表示同一类比较中的能量线，交点限定作用量的积分区间。有限高度的陡峭势垒还会产生随能量变化的反射相位，需由波函数与导数的连续条件求出；它通常介于软转点近似与无限硬壁极限之间。

还可把相位条件写成$J(E)=h[n+\eta(E)]$。对量子数求导，
\[
 \left[\frac{\dd J}{\dd E}-h\frac{\dd\eta}{\dd E}\right]\frac{\dd E}{\dd n}=h.
\]
若$\eta$为常数且半经典近似成立，便有$\Delta E\simeq h/T$，其中$T=\dd J/\dd E$是经典周期。若边界相位显著依赖能量，能级间隔还含$\dd\eta/\dd E$修正。无限深方阱可直接解边界方程，重力反弹低能级则应对照艾里函数零点；使用作用量时须区分精确边界解与渐近量子条件。

\section{光子统计与辐射能量尺度}
普朗克式中的$n$可以看作一个模态的光子占据数。黑体与腔壁可交换、产生和吸收光子，热平衡不限制光子总数，因而没有独立的光子数化学势。单模概率$p_n=(1-q)q^n$，其中$q=\ee^{-h\nu/(\kb T)}$。不同模态的占据独立时，总能量平均是各模能量平均之和；这不等于把每个频率区间强行塞入一个光子。

\begin{example}\label{ex:m4photonnumber}
求温度$T$的黑体辐射光子数密度$n_\gamma$和平均光子能量。可把$\zeta(3)=\sum_{j=1}^\infty j^{-3}$保留为常数。
\Diagram{fig:m4numberq}{能量谱除以单光子能量后才得到光子数谱}{\draw[curve](0,0)rectangle(2.3,1.1);\node at(1.15,.55){$u_\nu\dd\nu$};\draw[->](2.5,.55)--(3.5,.55)node[midway,above]{除$h\nu$};\node[right]at(3.7,.55){$\dd n_\gamma$};}
\end{example}
\begin{solution}
图\ref{fig:m4numberq}中，每个频率为$\nu$的光子能量是$h\nu$。把单位体积、单位频率能量密度除以它，再对频率积分：
\[
 n_\gamma=\int_0^\infty\frac{u_\nu}{h\nu}\dd\nu
 =\frac{8\pi\kb^3T^3}{h^3c^3}
 \int_0^\infty\frac{x^2}{\ee^x-1}\dd x,
 \quad x=\frac{h\nu}{\kb T}.
\]
展开$(\ee^x-1)^{-1}=\sum_{j\geq1}\ee^{-jx}$，各项非负，允许逐项积分。令$y=jx$、$\dd x=\dd y/j$，则
\[
 \int_0^\infty x^2\ee^{-jx}\dd x
 =j^{-3}\int_0^\infty y^2\ee^{-y}\dd y=\frac2{j^3}.
\]
最后一个积分的两次分部积分与端点计算已在例\ref{ex:m4stefan}给出。
\Diagram{fig:m4numbera}{光子数随$T^3$，总能量随$T^4$，平均能量随$T$}{\draw[axis](0,0)--(4.6,0)node[right]{$T$};\draw[axis](0,0)--(0,3)node[above]{$\langle E\rangle$};\draw[curve](0,0)--(4,2.5);\node[above]at(2.2,2.7){斜率$\pi^4\kb/[30\zeta(3)]$};}
积分因此为$2\zeta(3)$，得到
\[
 n_\gamma=\frac{16\pi\zeta(3)\kb^3T^3}{h^3c^3},\qquad
 \langle E\rangle=\frac{u}{n_\gamma}
 =\frac{\pi^4}{30\zeta(3)}\kb T\simeq2.701\kb T.
\]
图\ref{fig:m4numbera}给出按光子数加权的平均能量。相应概率密度为$p_N(\nu)=u_\nu/(h\nu n_\gamma)$，满足$\int p_N\dd\nu=1$，故$\langle E\rangle=\int h\nu p_N\dd\nu$。能量谱峰对应的$h\nu_{\max}\simeq2.821\kb T$则由谱函数的驻值确定。若先将能量谱归一化再平均$h\nu$，使用的已是能量权重，所得数值也会改变。
\end{solution}

\begin{example}\label{ex:m4fluctuation}
对单个频率为$\nu$的热平衡模态，求占据数方差和能量方差。讨论低频和高频极限。
\Diagram{fig:m4fluctq}{单模各占据数的几何概率分布}{\draw[axis](0,0)--(5.4,0)node[right]{$n$};\draw[axis](0,0)--(0,3)node[above]{$p_n$};\foreach\x in{0,1,2,3,4}{\draw[very thick](\x,0)--(\x,{2.5*.6^\x});}\node[below]at(0,0){$0$};}
\end{example}
\begin{solution}
对图\ref{fig:m4fluctq}的单模态，令$q=\ee^{-h\nu/(\kb T)}$，有$p_n=(1-q)q^n$。从几何级数$\sum_{n\geq0}q^n=(1-q)^{-1}$开始，分别施加算符$q\dd/\dd q$，
\[
 \sum nq^n=q\frac{\dd}{\dd q}(1-q)^{-1}
 =q[-(1-q)^{-2}(-1)]=\frac q{(1-q)^2},
\]
\[
 \sum n^2q^n=q\frac{\dd}{\dd q}[q(1-q)^{-2}]
 =q[(1-q)^{-2}+2q(1-q)^{-3}]
 =\frac{q(1+q)}{(1-q)^3}.
\]
乘以$(1-q)$即得$\bar n=q/(1-q)$、$\langle n^2\rangle=q(1+q)/(1-q)^2$。
\Diagram{fig:m4flucta}{方差由粒子项和二次涨落项组成}{\draw[axis](0,0)--(4.3,0)node[right]{$\bar n$};\draw[axis](0,0)--(0,3)node[above]{$\operatorname{Var}n$};\draw[curve,domain=0:1.4,samples=60]plot({2.3*\x},{.8*(\x+\x*\x)});\draw[guide](0,0)--(3.3,1.15);}
按方差定义相减，
\[
 \operatorname{Var}n=\frac{q(1+q)-q^2}{(1-q)^2}
 =\frac q{(1-q)^2}=\bar n(1+\bar n),
\]
\[
 \operatorname{Var}E=(h\nu)^2\bar n+(h\nu\bar n)^2.
\]
图\ref{fig:m4flucta}的两项在不同极限主导：高频$h\nu\gg\kb T$时$\bar n\ll1$，方差近于$\bar n$，但相对涨落$\sqrt{\operatorname{Var}n}/\bar n$很大；低频$\bar n\simeq\kb T/(h\nu)\gg1$，能量方差趋$(\kb T)^2$，单模态的相对能量涨落仍约为1。若测量包含$M$个独立、同分布模态，总平均能量乘$M$、方差乘$M$，相对涨落才减小为单模态值的$1/\sqrt M$。低频极限的平均模能为$\kb T$，单模测量仍有同量级的能量涨落。
\end{solution}

\begin{example}\label{ex:m4mixedcavity}
小腔中的某些方向由温度$T_1$的黑面照亮，其他方向由温度$T_2$的黑面照亮，立体角份额为$f,1-f$，腔内光不再充分热化。求角平均谱，判断它能否在所有频率上等价于单一温度的黑体。
\Diagram{fig:m4mixedq}{不同温度的方向混合一般不是单温度平衡谱}{\draw[curve](0,0)circle(1.6);\draw[curve](0,0)--(1.6,0);\draw[curve](0,0)--(-.8,1.386);\node at(.1,.8){$T_1$};\node at(0,-.7){$T_2$};\fill(0,0)circle(1.5pt);}
\end{example}
\begin{solution}
图\ref{fig:m4mixedq}中的两个方向集合互不重叠、各自的辐射互不相干，故角平均谱是强度密度之和：$\bar u_\nu=fu_\nu(T_1)+(1-f)u_\nu(T_2)$。按全频积分，
\[
 \bar u=a[fT_1^4+(1-f)T_2^4],\qquad
 T_E=[fT_1^4+(1-f)T_2^4]^{1/4},
\]
其中$a=4\sigma/c$。$T_E$只使某个单温黑体具有同一总能量，并未保证逐频相同。
\Diagram{fig:m4mixeda}{混合谱与等总能量单温谱；采用相同纵轴归一化}{
\draw[axis](0,0)--(4.8,0)node[right]{$h\nu/(\kb T_2)$};
\draw[axis](0,0)--(0,3)node[above]{谱密度};
\draw[curve]plot coordinates{(0.000000,0.0000000) (0.014400,0.0008289) (0.028800,0.0032713) (0.043200,0.0072621) (0.057600,0.0127375) (0.072000,0.0196350) (0.086400,0.0278933) (0.100800,0.0374527) (0.115200,0.0482547) (0.129600,0.0602420) (0.144000,0.0733589) (0.158400,0.0875507) (0.172800,0.1027643) (0.187200,0.1189476) (0.201600,0.1360501) (0.216000,0.1540223) (0.230400,0.1728162) (0.244800,0.1923848) (0.259200,0.2126826) (0.273600,0.2336651) (0.288000,0.2552893) (0.302400,0.2775131) (0.316800,0.3002959) (0.331200,0.3235979) (0.345600,0.3473808) (0.360000,0.3716074) (0.374400,0.3962414) (0.388800,0.4212479) (0.403200,0.4465930) (0.417600,0.4722437) (0.432000,0.4981683) (0.446400,0.5243361) (0.460800,0.5507174) (0.475200,0.5772835) (0.489600,0.6040067) (0.504000,0.6308602) (0.518400,0.6578183) (0.532800,0.6848562) (0.547200,0.7119499) (0.561600,0.7390765) (0.576000,0.7662137) (0.590400,0.7933404) (0.604800,0.8204360) (0.619200,0.8474809) (0.633600,0.8744564) (0.648000,0.9013443) (0.662400,0.9281274) (0.676800,0.9547892) (0.691200,0.9813139) (0.705600,1.0076863) (0.720000,1.0338921) (0.734400,1.0599175) (0.748800,1.0857493) (0.763200,1.1113752) (0.777600,1.1367833) (0.792000,1.1619623) (0.806400,1.1869015) (0.820800,1.2115909) (0.835200,1.2360208) (0.849600,1.2601823) (0.864000,1.2840667) (0.878400,1.3076660) (0.892800,1.3309728) (0.907200,1.3539798) (0.921600,1.3766804) (0.936000,1.3990685) (0.950400,1.4211381) (0.964800,1.4428840) (0.979200,1.4643010) (0.993600,1.4853847) (1.008000,1.5061307) (1.022400,1.5265350) (1.036800,1.5465942) (1.051200,1.5663049) (1.065600,1.5856643) (1.080000,1.6046696) (1.094400,1.6233185) (1.108800,1.6416090) (1.123200,1.6595392) (1.137600,1.6771076) (1.152000,1.6943129) (1.166400,1.7111540) (1.180800,1.7276300) (1.195200,1.7437404) (1.209600,1.7594848) (1.224000,1.7748628) (1.238400,1.7898746) (1.252800,1.8045203) (1.267200,1.8188001) (1.281600,1.8327147) (1.296000,1.8462646) (1.310400,1.8594508) (1.324800,1.8722742) (1.339200,1.8847358) (1.353600,1.8968370) (1.368000,1.9085792) (1.382400,1.9199637) (1.396800,1.9309923) (1.411200,1.9416666) (1.425600,1.9519886) (1.440000,1.9619600) (1.454400,1.9715831) (1.468800,1.9808598) (1.483200,1.9897924) (1.497600,1.9983832) (1.512000,2.0066345) (1.526400,2.0145489) (1.540800,2.0221287) (1.555200,2.0293767) (1.569600,2.0362954) (1.584000,2.0428876) (1.598400,2.0491560) (1.612800,2.0551034) (1.627200,2.0607328) (1.641600,2.0660469) (1.656000,2.0710489) (1.670400,2.0757416) (1.684800,2.0801282) (1.699200,2.0842116) (1.713600,2.0879951) (1.728000,2.0914818) (1.742400,2.0946748) (1.756800,2.0975774) (1.771200,2.1001928) (1.785600,2.1025243) (1.800000,2.1045750) (1.814400,2.1063484) (1.828800,2.1078478) (1.843200,2.1090764) (1.857600,2.1100377) (1.872000,2.1107350) (1.886400,2.1111716) (1.900800,2.1113510) (1.915200,2.1112766) (1.929600,2.1109517) (1.944000,2.1103799) (1.958400,2.1095644) (1.972800,2.1085087) (1.987200,2.1072163) (2.001600,2.1056906) (2.016000,2.1039350) (2.030400,2.1019529) (2.044800,2.0997477) (2.059200,2.0973230) (2.073600,2.0946820) (2.088000,2.0918283) (2.102400,2.0887651) (2.116800,2.0854961) (2.131200,2.0820244) (2.145600,2.0783536) (2.160000,2.0744870) (2.174400,2.0704280) (2.188800,2.0661799) (2.203200,2.0617461) (2.217600,2.0571299) (2.232000,2.0523347) (2.246400,2.0473638) (2.260800,2.0422205) (2.275200,2.0369080) (2.289600,2.0314297) (2.304000,2.0257888) (2.318400,2.0199885) (2.332800,2.0140321) (2.347200,2.0079227) (2.361600,2.0016636) (2.376000,1.9952580) (2.390400,1.9887089) (2.404800,1.9820195) (2.419200,1.9751930) (2.433600,1.9682324) (2.448000,1.9611408) (2.462400,1.9539212) (2.476800,1.9465767) (2.491200,1.9391103) (2.505600,1.9315250) (2.520000,1.9238237) (2.534400,1.9160093) (2.548800,1.9080849) (2.563200,1.9000533) (2.577600,1.8919172) (2.592000,1.8836797) (2.606400,1.8753436) (2.620800,1.8669115) (2.635200,1.8583863) (2.649600,1.8497707) (2.664000,1.8410675) (2.678400,1.8322793) (2.692800,1.8234088) (2.707200,1.8144586) (2.721600,1.8054314) (2.736000,1.7963297) (2.750400,1.7871562) (2.764800,1.7779132) (2.779200,1.7686034) (2.793600,1.7592292) (2.808000,1.7497931) (2.822400,1.7402974) (2.836800,1.7307446) (2.851200,1.7211371) (2.865600,1.7114771) (2.880000,1.7017670) (2.894400,1.6920091) (2.908800,1.6822056) (2.923200,1.6723587) (2.937600,1.6624707) (2.952000,1.6525437) (2.966400,1.6425798) (2.980800,1.6325812) (2.995200,1.6225499) (3.009600,1.6124880) (3.024000,1.6023975) (3.038400,1.5922803) (3.052800,1.5821386) (3.067200,1.5719742) (3.081600,1.5617889) (3.096000,1.5515848) (3.110400,1.5413636) (3.124800,1.5311272) (3.139200,1.5208773) (3.153600,1.5106158) (3.168000,1.5003444) (3.182400,1.4900647) (3.196800,1.4797785) (3.211200,1.4694874) (3.225600,1.4591931) (3.240000,1.4488971) (3.254400,1.4386010) (3.268800,1.4283063) (3.283200,1.4180147) (3.297600,1.4077275) (3.312000,1.3974462) (3.326400,1.3871723) (3.340800,1.3769072) (3.355200,1.3666522) (3.369600,1.3564088) (3.384000,1.3461781) (3.398400,1.3359617) (3.412800,1.3257606) (3.427200,1.3155762) (3.441600,1.3054097) (3.456000,1.2952623) (3.470400,1.2851352) (3.484800,1.2750295) (3.499200,1.2649463) (3.513600,1.2548868) (3.528000,1.2448520) (3.542400,1.2348430) (3.556800,1.2248608) (3.571200,1.2149064) (3.585600,1.2049808) (3.600000,1.1950850) (3.614400,1.1852198) (3.628800,1.1753862) (3.643200,1.1655851) (3.657600,1.1558174) (3.672000,1.1460838) (3.686400,1.1363853) (3.700800,1.1267226) (3.715200,1.1170965) (3.729600,1.1075077) (3.744000,1.0979571) (3.758400,1.0884452) (3.772800,1.0789729) (3.787200,1.0695407) (3.801600,1.0601494) (3.816000,1.0507995) (3.830400,1.0414917) (3.844800,1.0322266) (3.859200,1.0230048) (3.873600,1.0138267) (3.888000,1.0046931) (3.902400,0.9956043) (3.916800,0.9865609) (3.931200,0.9775634) (3.945600,0.9686122) (3.960000,0.9597078) (3.974400,0.9508506) (3.988800,0.9420410) (4.003200,0.9332795) (4.017600,0.9245664) (4.032000,0.9159021) (4.046400,0.9072870) (4.060800,0.8987213) (4.075200,0.8902054) (4.089600,0.8817396) (4.104000,0.8733242) (4.118400,0.8649595) (4.132800,0.8566457) (4.147200,0.8483831) (4.161600,0.8401719) (4.176000,0.8320124) (4.190400,0.8239046) (4.204800,0.8158489) (4.219200,0.8078454) (4.233600,0.7998943) (4.248000,0.7919957) (4.262400,0.7841498) (4.276800,0.7763567) (4.291200,0.7686165) (4.305600,0.7609294) (4.320000,0.7532953)};
\draw[guide]plot coordinates{(0.000000,0.0000000) (0.014400,0.0009450) (0.028800,0.0037358) (0.043200,0.0083068) (0.057600,0.0145934) (0.072000,0.0225320) (0.086400,0.0320602) (0.100800,0.0431165) (0.115200,0.0556406) (0.129600,0.0695730) (0.144000,0.0848556) (0.158400,0.1014309) (0.172800,0.1192429) (0.187200,0.1382362) (0.201600,0.1583567) (0.216000,0.1795513) (0.230400,0.2017678) (0.244800,0.2249551) (0.259200,0.2490631) (0.273600,0.2740428) (0.288000,0.2998459) (0.302400,0.3264256) (0.316800,0.3537356) (0.331200,0.3817309) (0.345600,0.4103673) (0.360000,0.4396016) (0.374400,0.4693918) (0.388800,0.4996966) (0.403200,0.5304757) (0.417600,0.5616899) (0.432000,0.5933006) (0.446400,0.6252706) (0.460800,0.6575634) (0.475200,0.6901432) (0.489600,0.7229755) (0.504000,0.7560264) (0.518400,0.7892631) (0.532800,0.8226536) (0.547200,0.8561668) (0.561600,0.8897724) (0.576000,0.9234410) (0.590400,0.9571441) (0.604800,0.9908540) (0.619200,1.0245439) (0.633600,1.0581875) (0.648000,1.0917598) (0.662400,1.1252363) (0.676800,1.1585933) (0.691200,1.1918080) (0.705600,1.2248582) (0.720000,1.2577227) (0.734400,1.2903808) (0.748800,1.3228128) (0.763200,1.3549994) (0.777600,1.3869223) (0.792000,1.4185638) (0.806400,1.4499069) (0.820800,1.4809353) (0.835200,1.5116333) (0.849600,1.5419860) (0.864000,1.5719790) (0.878400,1.6015987) (0.892800,1.6308321) (0.907200,1.6596666) (0.921600,1.6880904) (0.936000,1.7160924) (0.950400,1.7436618) (0.964800,1.7707887) (0.979200,1.7974635) (0.993600,1.8236772) (1.008000,1.8494215) (1.022400,1.8746884) (1.036800,1.8994705) (1.051200,1.9237611) (1.065600,1.9475538) (1.080000,1.9708426) (1.094400,1.9936222) (1.108800,2.0158876) (1.123200,2.0376344) (1.137600,2.0588585) (1.152000,2.0795562) (1.166400,2.0997244) (1.180800,2.1193603) (1.195200,2.1384615) (1.209600,2.1570259) (1.224000,2.1750521) (1.238400,2.1925387) (1.252800,2.2094848) (1.267200,2.2258900) (1.281600,2.2417538) (1.296000,2.2570766) (1.310400,2.2718587) (1.324800,2.2861009) (1.339200,2.2998042) (1.353600,2.3129699) (1.368000,2.3255998) (1.382400,2.3376955) (1.396800,2.3492594) (1.411200,2.3602937) (1.425600,2.3708012) (1.440000,2.3807847) (1.454400,2.3902473) (1.468800,2.3991923) (1.483200,2.4076233) (1.497600,2.4155440) (1.512000,2.4229584) (1.526400,2.4298704) (1.540800,2.4362845) (1.555200,2.4422051) (1.569600,2.4476368) (1.584000,2.4525843) (1.598400,2.4570527) (1.612800,2.4610469) (1.627200,2.4645721) (1.641600,2.4676338) (1.656000,2.4702372) (1.670400,2.4723880) (1.684800,2.4740919) (1.699200,2.4753545) (1.713600,2.4761817) (1.728000,2.4765795) (1.742400,2.4765539) (1.756800,2.4761110) (1.771200,2.4752570) (1.785600,2.4739980) (1.800000,2.4723404) (1.814400,2.4702905) (1.828800,2.4678547) (1.843200,2.4650395) (1.857600,2.4618512) (1.872000,2.4582965) (1.886400,2.4543818) (1.900800,2.4501138) (1.915200,2.4454989) (1.929600,2.4405438) (1.944000,2.4352551) (1.958400,2.4296395) (1.972800,2.4237034) (1.987200,2.4174536) (2.001600,2.4108967) (2.016000,2.4040393) (2.030400,2.3968879) (2.044800,2.3894491) (2.059200,2.3817296) (2.073600,2.3737358) (2.088000,2.3654744) (2.102400,2.3569517) (2.116800,2.3481742) (2.131200,2.3391484) (2.145600,2.3298807) (2.160000,2.3203775) (2.174400,2.3106450) (2.188800,2.3006896) (2.203200,2.2905175) (2.217600,2.2801349) (2.232000,2.2695479) (2.246400,2.2587626) (2.260800,2.2477851) (2.275200,2.2366214) (2.289600,2.2252774) (2.304000,2.2137590) (2.318400,2.2020720) (2.332800,2.1902221) (2.347200,2.1782152) (2.361600,2.1660568) (2.376000,2.1537524) (2.390400,2.1413077) (2.404800,2.1287281) (2.419200,2.1160189) (2.433600,2.1031855) (2.448000,2.0902331) (2.462400,2.0771668) (2.476800,2.0639919) (2.491200,2.0507134) (2.505600,2.0373361) (2.520000,2.0238651) (2.534400,2.0103051) (2.548800,1.9966609) (2.563200,1.9829372) (2.577600,1.9691386) (2.592000,1.9552697) (2.606400,1.9413348) (2.620800,1.9273385) (2.635200,1.9132849) (2.649600,1.8991785) (2.664000,1.8850232) (2.678400,1.8708233) (2.692800,1.8565828) (2.707200,1.8423055) (2.721600,1.8279955) (2.736000,1.8136565) (2.750400,1.7992922) (2.764800,1.7849063) (2.779200,1.7705024) (2.793600,1.7560840) (2.808000,1.7416546) (2.822400,1.7272175) (2.836800,1.7127761) (2.851200,1.6983335) (2.865600,1.6838929) (2.880000,1.6694575) (2.894400,1.6550302) (2.908800,1.6406141) (2.923200,1.6262119) (2.937600,1.6118266) (2.952000,1.5974609) (2.966400,1.5831174) (2.980800,1.5687989) (2.995200,1.5545078) (3.009600,1.5402466) (3.024000,1.5260178) (3.038400,1.5118238) (3.052800,1.4976668) (3.067200,1.4835491) (3.081600,1.4694729) (3.096000,1.4554402) (3.110400,1.4414532) (3.124800,1.4275139) (3.139200,1.4136242) (3.153600,1.3997860) (3.168000,1.3860011) (3.182400,1.3722713) (3.196800,1.3585982) (3.211200,1.3449837) (3.225600,1.3314293) (3.240000,1.3179365) (3.254400,1.3045068) (3.268800,1.2911418) (3.283200,1.2778427) (3.297600,1.2646110) (3.312000,1.2514480) (3.326400,1.2383548) (3.340800,1.2253328) (3.355200,1.2123831) (3.369600,1.1995068) (3.384000,1.1867050) (3.398400,1.1739786) (3.412800,1.1613287) (3.427200,1.1487563) (3.441600,1.1362621) (3.456000,1.1238471) (3.470400,1.1115121) (3.484800,1.0992578) (3.499200,1.0870851) (3.513600,1.0749945) (3.528000,1.0629867) (3.542400,1.0510624) (3.556800,1.0392222) (3.571200,1.0274665) (3.585600,1.0157960) (3.600000,1.0042110) (3.614400,0.9927121) (3.628800,0.9812997) (3.643200,0.9699741) (3.657600,0.9587356) (3.672000,0.9475847) (3.686400,0.9365216) (3.700800,0.9255466) (3.715200,0.9146598) (3.729600,0.9038615) (3.744000,0.8931520) (3.758400,0.8825312) (3.772800,0.8719994) (3.787200,0.8615567) (3.801600,0.8512031) (3.816000,0.8409386) (3.830400,0.8307634) (3.844800,0.8206774) (3.859200,0.8106806) (3.873600,0.8007729) (3.888000,0.7909543) (3.902400,0.7812247) (3.916800,0.7715841) (3.931200,0.7620322) (3.945600,0.7525690) (3.960000,0.7431942) (3.974400,0.7339078) (3.988800,0.7247095) (4.003200,0.7155991) (4.017600,0.7065763) (4.032000,0.6976410) (4.046400,0.6887929) (4.060800,0.6800317) (4.075200,0.6713571) (4.089600,0.6627688) (4.104000,0.6542665) (4.118400,0.6458498) (4.132800,0.6375185) (4.147200,0.6292721) (4.161600,0.6211104) (4.176000,0.6130329) (4.190400,0.6050392) (4.204800,0.5971289) (4.219200,0.5893017) (4.233600,0.5815571) (4.248000,0.5738947) (4.262400,0.5663140) (4.276800,0.5588146) (4.291200,0.5513961) (4.305600,0.5440579) (4.320000,0.5367997)};
\node[above]at(3.7,2.5){实线：混合谱};\node[above]at(3.7,2.1){虚线：$T_E$单温谱};
}
在低频端，普朗克谱$u_\nu\simeq8\pi\kb T\nu^2/c^3$，所以混合谱对应的低频温度为$T_L=fT_1+(1-f)T_2$。函数$T^4$在正温区严格凸，故对$0<f<1$、$T_1\ne T_2$有
\[
 T_E^4=fT_1^4+(1-f)T_2^4>[fT_1+(1-f)T_2]^4=T_L^4.
\]
因此同一温度不能同时匹配低频系数与总能量。图\ref{fig:m4mixeda}采用$T_1=2T_2$、$f=1/2$绘出两条普朗克组合曲线；$T_L=1.5T_2$、$T_E=8.5^{1/4}T_2\simeq1.7075T_2$。

高频端若$T_1>T_2$，较热分量的指数衰减更慢，也不会变成$T_E$的单一指数。只有两温相同或某份额为零等退化情形才得到全频单温黑体。若腔壁吸收、再发射足够充分，原来分方向保留的谱会被热化，本题不再适用。
\end{solution}

\section{能级间隔、周期与边界相位}
对一维哈密顿量$H=p^2/(2m)+V(x)$，在允许区$\partial p/\partial E=m/p=1/\dot x$。对闭合轨道作用量求能量导数，转点处$p=0$使端点贡献为零，得到
\begin{equation}\label{eq:m4period}
 \frac{\dd J}{\dd E}=\oint\frac{\dd x}{\dot x}=T(E).
\end{equation}
相邻半经典能级$\Delta J=h$，所以大量子数时$\Delta E\simeq h/T(E)$，即跃迁频率趋于经典周期频率。周期变化与能级疏密由同一个相空间面积导数决定。

\begin{example}\label{ex:m4period}
由作用量求重力反弹粒子的经典周期，并用它估计高能级间距，核对直接对能级公式求导的结果。
\Diagram{fig:m4periodq}{粒子由地面反弹到$x_t$后返回}{\draw[axis](0,0)--(0,3.7)node[above]{$x$};\draw[very thick](-1,0)--(1,0);\draw[->](-.4,.3)--(-.4,2.7);\draw[->](.4,2.7)--(.4,.3);\draw[guide](-1,2.7)--(1,2.7);\node[right]at(1,2.7){$x_t=E/(mg)$};}
\end{example}
\begin{solution}
图\ref{fig:m4periodq}的重力反弹轨道作用量为$J(E)=4\sqrt2E^{3/2}/(3g\sqrt m)$，对能量求导得到
\[
 T(E)=\frac{\dd J}{\dd E}
 =\frac{4\sqrt2}{3g\sqrt m}\frac32E^{1/2}
 =\frac{2\sqrt2}{g\sqrt m}\sqrt E
 =\frac{2v_0}g,\qquad v_0=\sqrt{2E/m}.
\]
这与粒子从硬壁起跳、到最高点、再回到硬壁的时间一致。硬壁瞬时反弹不增加有限飞行时间。
\Diagram{fig:m4perioda}{周期随能量增大，半经典相邻能级间距减小}{\draw[axis](0,0)--(4.5,0)node[right]{$E$};\draw[axis](0,0)--(0,3)node[above]{$T$};\draw[curve,domain=0:4,samples=70]plot(\x,{1.2*sqrt(\x)});}
图\ref{fig:m4perioda}中周期随能量增大。半经典条件$J(E_n)=h(n+3/4)$相邻相减，若$n$足够大，可将$J(E_{n+1})-J(E_n)$展开为$T(E_n)\Delta E+O(\Delta E^2)$，所以
\[
 \Delta E\simeq\frac h{T(E)}=\frac{hg\sqrt m}{2\sqrt{2E}}.
\]
直接从$E_n=C(n+3/4)^{2/3}$求导，
\[
 \frac{\dd E}{\dd n}=\frac23C(n+3/4)^{-1/3}
 =\frac{2C^{3/2}}{3\sqrt E}.
\]
再用$C^{3/2}=3hg\sqrt m/(4\sqrt2)$，恰得同一表达式。这里$h/T$估计相邻高能级间距。最低几个能级应保留有限差分，并与完整能级式或精确艾里函数根比较。经典周期越长，对应的相位随能量变化越快，达到下一个量子相位所需的能量增量越小。
\end{solution}

\begin{example}\label{ex:m4correspondence}
非相对论玻尔氢样原子中，比较$n\to n-1$跃迁频率和第$n$条圆轨道的转动频率，求大$n$极限下两者之比。
\Diagram{fig:m4correspondenceq}{相邻高量子数轨道间的跃迁}{\draw[curve](0,0)circle(1.6);\draw[curve](0,0)circle(1.2);\draw[->](1.6,0)--(1.2,0);\node[right]at(1.7,.3){$n$};\node[right]at(1.1,-.6){$n-1$};\fill(0,0)circle(2pt);}
\end{example}
\begin{solution}
令$C=mc^2(Z\alpha)^2/2$，玻尔能级为$E_n=-C/n^2$。图\ref{fig:m4correspondenceq}中$n\to n-1$释放的光子能量为
\[
 h\nu_{n,n-1}=C\left[\frac1{(n-1)^2}-\frac1{n^2}\right]
 =C\frac{2n-1}{n^2(n-1)^2},\quad n\geq2.
\]
第$n$条圆轨道的$v_n=Z\alpha c/n$、$r_n=\hbar n^2/(mZ\alpha c)$，所以转动频率为
\[
 \nu_{\rm orb}=\frac{v_n}{2\pi r_n}=\frac{2C}{hn^3}.
\]
\Diagram{fig:m4correspondencea}{高能级局部斜率给出跃迁频率的领先项}{\draw[axis](0,0)--(4.8,0)node[right]{$n$};\draw[axis](0,-2.7)--(0,.5)node[above]{$E_n$};\draw[curve,domain=1:4.2,samples=70]plot(\x,{-2.3/(\x*\x)});\foreach\x in{1,2,3,4}{\fill(\x,{-2.3/(\x*\x)})circle(1.5pt);}}
两者之比可不经近似直接写为
\[
 \frac{\nu_{n,n-1}}{\nu_{\rm orb}}
 =\frac{n(2n-1)}{2(n-1)^2}
 =1+\frac{3}{2n}+\frac{2}{n^2}+O(n^{-3}).
\]
图\ref{fig:m4correspondencea}中的两条频率在大$n$处相近，极限比值为1。展开时用$(1-1/n)^{-2}=1+2/n+3/n^2+\cdots$，再乘$1-1/(2n)$，才能保留正确的一阶修正。对固定整数$k\ll n$，$n\to n-k$的跃迁频率近于$k\nu_{\rm orb}$，对应经典周期运动的第$k$次谐波；这不等于有限$n$的每条光谱线都能视作电子沿单一圆轨道连续辐射。
\end{solution}

\begin{example}\label{ex:m4finitebarrier}
粒子在$0<x<L$势能为零，两侧势垒均为$V_0>E$。求束缚态的相位量子条件，并说明无限深方阱是怎样的极限。只推导允许能量的方程，不要求逐根求值。
\Diagram{fig:m4finiteq}{有限深方阱在壁外有指数尾部}{\draw[axis](-1.6,0)--(5,0)node[right]{$x$};\draw[curve](-1.3,2.4)--(0,2.4)--(0,0)--(3,0)--(3,2.4)--(4.5,2.4);\draw[guide](-1.3,1)--(4.5,1);\node[right]at(4.5,1){$E$};\node[left]at(-1.3,2.4){$V_0$};}
\end{example}
\begin{solution}
图\ref{fig:m4finiteq}中令$k=\sqrt{2mE}/\hbar$、$\kappa=\sqrt{2m(V_0-E)}/\hbar$。左侧选$A\ee^{\kappa x}$，排除$x\to-\infty$增长的支；阱内取$B\sin(kx+\phi)$。在$x=0$连续匹配波函数及导数：
\[
 A=B\sin\phi,\qquad \kappa A=kB\cos\phi.
\]
两式相除得$k\cot\phi=\kappa$，可选$0<\phi<\pi/2$，故$\phi=\arctan(k/\kappa)$。
\Diagram{fig:m4finitea}{有限势阱边界处波函数非零，壁外有连续衰减尾部}{\draw[axis](-1.3,0)--(4.6,0)node[right]{$x$};\draw[guide](0,-.4)--(0,2.3);\draw[guide](3,-.4)--(3,2.3);\draw[curve,domain=-1:0,samples=35]plot(\x,{1.8*cos(67.5)*exp(1.896119*\x)});\draw[curve,domain=0:3,samples=80]plot(\x,{1.8*cos(45*(\x-1.5))});\draw[curve,domain=3:4,samples=35]plot(\x,{1.8*cos(67.5)*exp(-1.896119*(\x-3))});}
图\ref{fig:m4finitea}的右尾取$D\ee^{-\kappa(x-L)}$。在$x=L$，$D=B\sin(kL+\phi)$、$-\kappa D=kB\cos(kL+\phi)$，所以
\[
 \cot(kL+\phi)=-\cot\phi,
 \quad kL+\phi=n\pi-\phi,
\]
得到允许能量方程
\[
 kL+2\arctan(k/\kappa)=n\pi,
 \qquad 0<E<V_0,\quad n=1,2,\ldots.
\]
求根须同时代入$k(E)$和$\kappa(E)$，不能把它们视为独立变量。左端随$E$从0单调增至$L\sqrt{2mV_0}/\hbar+\pi$的开端点；只有严格低于这个上限的$n\pi$才有束缚态根，因此根的数目有限。$V_0\to\infty$时$\kappa\to\infty$、附加相位趋零，恢复无限深阱。有限势垒允许波函数渗出，等效约束较弱，同一编号的能量低于无限深阱相应能量。
\end{solution}

\begin{example}\label{ex:m4triangular}
一维势$V(x)=F|x|$，$F>0$，无硬壁。求半经典能级；与只保留$x\geq0$并在原点设置硬壁的模型比较相位偏移和作用量。
\Diagram{fig:m4triq}{对称线性势具有两端软转点，原点不是硬壁}{\draw[axis](-2.7,0)--(2.9,0)node[right]{$x$};\draw[axis](0,0)--(0,3.2)node[above]{$V$};\draw[curve](-2.4,2.7)--(0,0)--(2.4,2.7);\draw[guide](-1.8,2.025)--(1.8,2.025);\node[right]at(1.8,2.025){$E$};}
\end{example}
\begin{solution}
图\ref{fig:m4triq}的转点为$\pm E/F$。完整作用量的四倍因子计入左右两半区及往返两支。取$q=E-Fx$，$\dd x=-\dd q/F$，有
\begin{align*}
 J&=4\int_0^{E/F}\sqrt{2m(E-Fx)}\dd x
 =-\frac{4\sqrt{2m}}F\int_E^0q^{1/2}\dd q\\
 &=\frac{4\sqrt{2m}}F\left[\frac23q^{3/2}\right]_0^E
 =\frac{8\sqrt{2m}}{3F}E^{3/2}.
\end{align*}
取两个软转点条件，
\[
 E_n\simeq\left[\frac{3Fh}{8\sqrt{2m}}(n+1/2)\right]^{2/3}.
\]
\Diagram{fig:m4tria}{对称势相轨道跨越原点；硬壁模型只保留右半边}{\draw[axis](-2.6,0)--(2.8,0)node[right]{$x$};\draw[axis](0,-2)--(0,2.2)node[above]{$p$};\draw[curve,domain=-2:2,samples=80]plot(\x,{1.6*sqrt(1-abs(\x)/2)});\draw[curve,domain=-2:2,samples=80]plot(\x,{-1.6*sqrt(1-abs(\x)/2)});}
若原点改硬壁，图\ref{fig:m4tria}只留右半区，相轨道面积减半，且条件变为$(n+3/4)h$，能级为$[3Fh(n+3/4)/(4\sqrt{2m})]^{2/3}$。这对应对称势的奇宇称能级子集的半经典极限，不能只更换偏移而不改积分区域。原点势的尖点不是反射硬壁，波函数可以穿过。
\end{solution}

\begin{example}\label{ex:m4photonrecoil}
自由原子初始静止，激发态与基态静质量分别为$M_i,M_f$。原子辐射一个光子退激，求光子能量和基态原子的反冲能，比较内部能差$\Delta=(M_i-M_f)c^2$。
\Diagram{fig:m4recoilq}{自发辐射必须连同原子反冲一起守恒}{\draw[curve](0,0)circle(.45);\node at(0,0){$M_i$};\draw[->](.7,0)--(2,0);\draw[->](2.6,.4)--(4.7,.4)node[above]{$E_\gamma$};\draw[->](2.4,-.4)--(1.2,-.4)node[below]{$M_f$反冲};}
\end{example}
\begin{solution}
在图\ref{fig:m4recoilq}的原子初始静止系，末光子沿$+x$，其能量为$E_\gamma$。基态原子反冲动量为$-E_\gamma/c$，总能量为$M_ic^2-E_\gamma$。对末原子写质壳条件，
\begin{align*}
 M_f^2c^4&=(M_ic^2-E_\gamma)^2-E_\gamma^2\\
 &=M_i^2c^4-2M_ic^2E_\gamma.
\end{align*}
解得$E_\gamma=(M_i^2-M_f^2)c^2/(2M_i)$。
\Diagram{fig:m4recoila}{内部能差包含光子能量和反冲动能}{\draw[curve](0,0)rectangle(5.4,1);\draw[curve](4.5,0)--(4.5,1);\node at(2.25,.5){$E_\gamma$};\node at(4.95,.5){$K_R$};\node[above]at(2.7,1){$\Delta$};}
令内部能差$\Delta=(M_i-M_f)c^2$，利用$M_f=M_i-\Delta/c^2$展开平方，得到
\[
 E_\gamma=\Delta\left(1-\frac{\Delta}{2M_ic^2}\right),\quad
 K_R=M_ic^2-E_\gamma-M_fc^2=\frac{\Delta^2}{2M_ic^2}.
\]
图\ref{fig:m4recoila}中的能量分配$E_\gamma+K_R=\Delta$是精确的两体运动学结果；$K_R=p_R^2/(2M_f)$仅是非相对论近似，不能先用它求取精确公式。末原子速率为
\[
 |V_R|=\frac{p_Rc^2}{E_R}
 =c\frac{M_i^2-M_f^2}{M_i^2+M_f^2}.
\]
当$\Delta\ll M_ic^2$时$|V_R|\simeq\Delta/(M_ic)$，反冲修正很小。若讨论反向过程，即静止基态原子吸收一个光子而整体升为激发态，所需能量为$(M_i^2-M_f^2)c^2/(2M_f)>\Delta$；发射线与自由吸收线因反冲而不重合。
\end{solution}

\section{有限束缚边界与势阱的时间改变}
\subsection{允许能量、归一化和边界外概率}
有限势阱的束缚态在禁阻区保留指数尾部。边界处波函数及其导数的连续性决定允许能级，完整波函数的归一化再决定阱内、阱外概率。这两步应使用同一套分段解。

取对称阱$|x|<a/2$内$V=0$，外侧$V=V_0$。束缚态满足$0<E<V_0$，定义$k=\sqrt{2mE}/\hbar$、$\kappa=\sqrt{2m(V_0-E)}/\hbar$。偶态写成
\[
 \psi(x)=\begin{cases}
 A\cos kx,&|x|<a/2,\\
 A\cos(ka/2)\ee^{-\kappa(|x|-a/2)},&|x|\geq a/2.
 \end{cases}
\]
在右边界令导数连续，得$k\tan(ka/2)=\kappa$。奇态相应给$-k\cot(ka/2)=\kappa$。令$z=ka/2$、$z_0=a\sqrt{2mV_0}/(2\hbar)$，则$\kappa a/2=\sqrt{z_0^2-z^2}$。应在每个正切或余切分支上寻找合法交点，排除$E=V_0$时不能归一化的阈值态。

偶态的归一化包含阱内余弦段和两个尾部。尾部取$s=|x|-a/2$，各自的上下限为$0,\infty$，于是
\begin{align*}
 A^{-2}&=\int_{-a/2}^{a/2}\cos^2kx\dd x
 +2\cos^2(ka/2)\int_0^\infty\ee^{-2\kappa s}\dd s\\
 &=\left[\frac x2+\frac{\sin2kx}{4k}\right]_{-a/2}^{a/2}
 +2\cos^2(ka/2)\left[-\frac{\ee^{-2\kappa s}}{2\kappa}\right]_0^\infty\\
 &=\frac a2+\frac{\sin ka}{2k}+\frac{\cos^2(ka/2)}\kappa.
\end{align*}
最后一项正是禁阻区概率的未归一化权重。$V_0$很大时$\kappa$很大，尾部减小；靠近上沿的束缚态$\kappa$很小，延伸长度$1/\kappa$可能很长，不能仅用阱内几何宽度估计波函数范围。

\begin{example}一个能精确计算尾部概率的有限阱
上述对称阱的基态满足$k=\kappa$。求所需$V_0$、基态能量、归一化常数及粒子处于阱外的概率。
\ExpansionDiagram{fig:dep-m4-well-q}{有限方势阱的宽度与基态能级，阱内势取零}{
\draw[->](-3,0)--(3.2,0)node[right]{$x$};\draw[thick](-3,2)--(-1,2)--(-1,0)--(1,0)--(1,2)--(3,2);\draw[dashed](-2.8,1)--(2.8,1)node[right]{$E$};\node[above]at(2,2){$V_0$};\node[below]at(-1,0){$-a/2$};\node[below]at(1,0){$a/2$};}
\end{example}
\begin{solution}
基态没有内部节点，$0<ka/2<\pi/2$。由$k\tan(ka/2)=\kappa=k$，得$ka/2=\pi/4$，所以$k=\kappa=\pi/(2a)$。于是
\[
 E=\frac{\pi^2\hbar^2}{8ma^2},\qquad
 V_0-E=E,\qquad V_0=\frac{\pi^2\hbar^2}{4ma^2}.
\]
图\ref{fig:dep-m4-well-q}中的能级恰在阱深一半，但不能由这个条件直接猜测阱外概率也为一半。代入归一化公式，$\sin ka=1$、$\cos^2(ka/2)=1/2$，得
\[
 A^{-2}=\frac a2+\frac a\pi+\frac a\pi
 =\frac{a(\pi+4)}{2\pi},\qquad
 A=\sqrt{\frac{2\pi}{a(\pi+4)}}.
\]
\ExpansionDiagram{fig:dep-m4-well-a}{波函数的余弦段与指数尾部在边界平滑接续}{
\draw[->](-3,0)--(3.2,0)node[right]{$x/(a/2)$};\draw[thick,domain=-1:1,samples=80]plot(\x,{1.8*cos(45*\x)});\draw[thick,domain=1:3,samples=60]plot(\x,{1.8/sqrt(2)*exp(-pi*(\x-1)/4)});\draw[thick,domain=-3:-1,samples=60]plot(\x,{1.8/sqrt(2)*exp(pi*(\x+1)/4)});\draw[dashed](-1,0)--(-1,2);\draw[dashed](1,0)--(1,2);\draw[->](-3,0)--(-3,2.4)node[above]{$\psi/A$};\draw(-3.06,1.8)--(-2.94,1.8);\node[left]at(-3,1.8){$1$};\foreach \x in{-1,0,1}{\draw(\x,-.04)--(\x,.04);\node[below]at(\x,0){$\x$};}}
由图\ref{fig:dep-m4-well-a}两侧尾部积分，
\[
 P_{\rm out}=A^2\frac{\cos^2(ka/2)}\kappa
 =\frac{2}{\pi+4}\simeq0.2800.
\]
所以约28\%概率位于经典禁阻区。这里“禁阻”指经典动能$E-V_0<0$，并不要求量子波函数严格为零。若只在阱内归一化，会把所有局部概率密度都高估。
\end{solution}

\subsection{突然改变与缓慢改变使用不同的接续条件}
势阱宽度缓慢改变时，满足绝热条件的非简并本征态沿同一能级分支演化；突然改变时，波函数来不及改变，应把原波函数投影到新本征态。两者都需要初态和改变方式，不能只因初末几何宽度相同便使用同一能量公式。

\begin{example}无限深阱突然加宽为两倍
粒子最初处于$0<x<L$无限深阱基态。右壁瞬间移到$2L$，左壁不动。求改变后各能级的概率、能量期望，并与缓慢加宽的结果比较。取瞬间改变为理想极限，外部势在新增区间中直接降为零。
\ExpansionDiagram{fig:dep-m4-sudden-q}{右壁移开的一瞬间，原波函数在新增区间仍为零}{
\draw[->](0,0)--(6,0)node[right]{$x$};\draw[very thick](.5,0)--(.5,2);\draw[dashed,very thick](3,0)--(3,2);\draw[very thick](5.5,0)--(5.5,2);\draw[thick,domain=.5:3,samples=70]plot(\x,{1.5*sin(180*(\x-.5)/2.5)});\draw[thick](3,0)--(5.5,0);\node[below]at(3,0){$L$};\node[below]at(5.5,0){$2L$};\draw[->](3,2.3)--(5.5,2.3);}
\end{example}
\begin{solution}
图\ref{fig:dep-m4-sudden-q}给定瞬间波函数：$\psi=\sqrt{2/L}\sin(\pi x/L)$于$0<x<L$，在$L<x<2L$为零。新阱归一化本征函数为$\phi_k=L^{-1/2}\sin(k\pi x/(2L))$，能级为$E_k'=k^2\pi^2\hbar^2/(8mL^2)$。系数
\[
 c_k=\frac{\sqrt2}{L}\int_0^L\sin\frac{\pi x}{L}
 \sin\frac{k\pi x}{2L}\dd x.
\]
对$k\ne2$，先使用$\sin a\sin b=[\cos(a-b)-\cos(a+b)]/2$，再分别积分：
\begin{align*}
 c_k&=\frac{\sqrt2}{2L}\int_0^L
 \left[\cos\frac{(2-k)\pi x}{2L}-\cos\frac{(2+k)\pi x}{2L}\right]\dd x\\
 &=\frac{\sqrt2}{\pi}
 \left[\frac{\sin((2-k)\pi x/(2L))}{2-k}
 -\frac{\sin((2+k)\pi x/(2L))}{2+k}\right]_0^L\\
 &=\frac{\sqrt2}{\pi}\sin\frac{k\pi}{2}
 \left(\frac1{2-k}+\frac1{2+k}\right).
\end{align*}
这里用了$\sin(\pi-k\pi/2)=\sin(k\pi/2)$、$\sin(\pi+k\pi/2)=-\sin(k\pi/2)$，下端点两项均为零。因此
\[
 c_k=\frac{4\sqrt2\sin(k\pi/2)}{\pi(4-k^2)}.
\]
$k=2$时应直接积分两个相同的正弦函数：
\[
 c_2=\frac{\sqrt2}{L}\int_0^L\sin^2\frac{\pi x}{L}\dd x
 =\frac{\sqrt2}{L}\left[\frac x2-\frac{L}{4\pi}\sin\frac{2\pi x}{L}\right]_0^L
 =\frac1{\sqrt2}.
\]
其余偶数$k$的系数为零，奇数$k$则
\[
 P_k=|c_k|^2=\frac{32}{\pi^2(4-k^2)^2},
 \qquad P_2=\frac12.
\]
\ExpansionDiagram{fig:dep-m4-sudden-a}{突然加宽后能量测量分布在多个新能级上}{
\draw[->](0,0)--(7,0)node[right]{$k$};\draw[->](0,0)--(0,3.5)node[above]{$P_k$};\draw[very thick](1,0)--(1,2.16);\draw[very thick](2,0)--(2,3);\draw[very thick](3,0)--(3,.78);\draw[very thick](5,0)--(5,.044);\draw[dashed](0,3)--(2,3);\node[left]at(0,3){$1/2$};\node[below]at(1,0){1};\node[below]at(2,0){2};\node[below]at(3,0){3};\node[below]at(5,0){5};}
图\ref{fig:dep-m4-sudden-a}以同一纵向比例画前几个概率，其中$P_1=32/(9\pi^2)\simeq0.3603$、$P_3=32/(25\pi^2)\simeq0.1297$。其总和为1，由新本征函数的完备性和初态归一化保证。

能量期望可直接用动能二次型，无需逐项硬求级数：
\begin{align*}
 \langle E'\rangle&=\frac{\hbar^2}{2m}\int_0^{2L}|\psi'(x)|^2\dd x
 =\frac{\pi^2\hbar^2}{mL^3}\int_0^L\cos^2\frac{\pi x}{L}\dd x\\
 &=\frac{\pi^2\hbar^2}{mL^3}
 \left[\frac x2+\frac{L}{4\pi}\sin\frac{2\pi x}{L}\right]_0^L\\
 &=\frac{\pi^2\hbar^2}{2mL^2}=E_{1,\rm old}.
\end{align*}
瞬间移去远侧势垒没有改变原波函数的动能期望；之后波包由多个新定态叠加演化。若缓慢加宽且保持基态分支，末态只有$k=1$，能量为原来的$1/4$，能量差由移动壁面的做功带走。突然加宽后的测量结果虽然包含低能级，也包含更高能级，不能只取最可能能级代替平均能量。
\end{solution}

\section{自由原子的吸收共振与反冲}
\subsection{内部能差不是所有过程中的同一个光子能量}
设自由原子基态、激发态静质量为$M,M_*$，内部能差$\Delta=(M_*-M)c^2>0$。激发态原子在静止时发射一个光子的能量已由两体守恒给出$E_{\rm em}=(M_*^2-M^2)c^2/(2M_*)$。反过来，基态静止原子吸收一个光子而成为同一激发态时，末态原子也要携带光子动量，吸收能量一般不同。

吸收过程初态总四动量平方为$M^2c^4+2Mc^2E_{\rm abs}$，末态为$M_*^2c^4$，所以
\[
 E_{\rm abs}=\frac{(M_*^2-M^2)c^2}{2M}
 =\Delta+\frac{\Delta^2}{2Mc^2}.
\]
吸收光子须同时供给内部激发能和反冲动能；发射时，内部能差则由光子与反冲原子分享。两题都把各自的初态原子设为静止，因此所选初末四动量没有互换，时间反演对称性并不要求这两个光子同频。

\begin{example}静止自由原子能否吸收同种静止原子发出的光
两种完全相同的自由原子最初分别处于激发态和基态，且都在实验室静止。第一原子辐射的单光子沿连线入射第二原子。只考虑同一窄能级跃迁，求发射能量与静止吸收共振能量的差，并估计为补偿该差所需的吸收原子速度。取$\Delta\ll Mc^2$，忽略自然线宽以外的其他展宽机制。
\ExpansionDiagram{fig:dep-m4-recoil-q}{发射原子与吸收原子各自需要分配反冲能}{
\fill(0,0)circle(3pt)node[below]{$M_*$};\draw[->,thick](.5,0)--(4.5,0)node[midway,above]{$E_{\rm em}$};\fill(5,0)circle(3pt)node[below]{$M$};\draw[->](0,-.6)--(-1,-.6)node[left]{发射反冲};\draw[->](5,-.6)--(6,-.6)node[right]{吸收反冲};}
\end{example}
\begin{solution}
图\ref{fig:dep-m4-recoil-q}中两原子的初态都静止，但两个末态都要反冲。能量差精确为
\begin{align*}
 E_{\rm abs}-E_{\rm em}
 &=\frac{(M_*^2-M^2)c^2}{2}
 \left(\frac1M-\frac1{M_*}\right)\\
 &=\frac{\Delta^2(M_*+M)}{2MM_*c^2}
 \simeq\frac{\Delta^2}{Mc^2}.
\end{align*}
这一正差约为两次单原子反冲能$E_R=\Delta^2/(2Mc^2)$之和。若线宽远小于$2E_R$，静止发射源的中心光子不足以达到静止吸收原子的共振能量，尽管两原子内部能级完全相同。
\ExpansionDiagram{fig:dep-m4-recoil-a}{发射与吸收的中心能量分处内部能差两侧}{
\draw[->](0,0)--(7,0)node[right]{光子能量};\draw[thick](2,-.4)--(2,1)node[above]{$E_{\rm em}$};\draw[dashed](3.5,-.4)--(3.5,1)node[above]{$\Delta$};\draw[thick](5,-.4)--(5,1)node[above]{$E_{\rm abs}$};\draw[<->](2,-.8)--(5,-.8)node[midway,below]{$\simeq2E_R$};}
若吸收原子以小速度$u$迎着入射光运动，在原子系测得光子能量约为$E_{\rm em}(1+u/c)$。共振要求$u/c\simeq(E_{\rm abs}-E_{\rm em})/\Delta\simeq\Delta/(Mc^2)$，即$u\simeq\Delta/(Mc)$。这次一阶多普勒蓝移恰好补偿图\ref{fig:dep-m4-recoil-a}中的能量缺口。

若气体处于温度$T$，沿光束方向的速度标准差为$\sqrt{\kb T/M}$，多普勒能量标准差约为$\sigma_E=\Delta\sqrt{\kb T/(Mc^2)}$。它与反冲缺口同量级时，$\kb T\sim\Delta^2/(Mc^2)$。计算完整吸收率还需要线形、速度分布和跃迁强度。原子若受外部结构约束，动量可由更大系统承担，能量和动量守恒也须对这个整体列式。
\end{solution}

\section{习题与解答}
\begin{exercise}\label{ex:m4tunnel}
质量$m$的粒子以能量$0<E<V_0$从左射向$0<x<a$的矩形势垒，势垒外$V=0$。求透射率，并给出厚势垒极限，解释它与束缚态边界的联系。
\Diagram{fig:m4tunnelq}{有限禁区两侧都是传播区}{\draw[axis](-1.5,0)--(5,0)node[right]{$x$};\draw[curve](-1.2,0)--(0,0)--(0,2.4)--(3,2.4)--(3,0)--(4.7,0);\draw[guide](-1.2,1)--(4.7,1);\node[right]at(4.7,1){$E$};\node[above]at(1.5,2.4){$V_0$};\node[below]at(3,0){$a$};}
\end{exercise}
\begin{solution}
图\ref{fig:m4tunnelq}取$k=\sqrt{2mE}/\hbar$、$\kappa=\sqrt{2m(V_0-E)}/\hbar$。三段波函数分别为$\ee^{ikx}+r\ee^{-ikx}$、$A\ee^{\kappa x}+B\ee^{-\kappa x}$、$t\ee^{ikx}$。右界面两连续条件给
\[
 A\ee^{\kappa a}=\frac{t\ee^{ika}}2(1+ik/\kappa),\quad
 B\ee^{-\kappa a}=\frac{t\ee^{ika}}2(1-ik/\kappa).
\]
代回左界面，先算出
\begin{align*}
 A+B&=t\ee^{ika}[\cosh\kappa a-i(k/\kappa)\sinh\kappa a],\\
 \kappa(A-B)&=t\ee^{ika}[-\kappa\sinh\kappa a+ik\cosh\kappa a].
\end{align*}
左界面条件为$1+r=A+B$、$ik(1-r)=\kappa(A-B)$。第二式除以$ik$，再与第一式相加，消去$r$，得
\[
 t\ee^{ika}=\left[\cosh\kappa a+
 i\frac{\kappa^2-k^2}{2k\kappa}\sinh\kappa a\right]^{-1}.
\]
\Diagram{fig:m4tunnela}{有限禁区需要两指数支，透射侧只保留向右传播波}{\draw[axis](-1.5,0)--(5,0)node[right]{$x$};\draw[guide](0,-.6)--(0,2.3);\draw[guide](3,-.6)--(3,2.3);\draw[curve,domain=0:3,samples=60]plot(\x,{1.8*exp(-.65*\x)});\draw[->](3.4,.6)--(4.7,.6)node[above]{透射};\node[above]at(1.5,2.4){振幅包络示意};}
两侧传播介质相同，$T=|t|^2$。分母的模平方为
\begin{align*}
 T^{-1}&=\cosh^2\kappa a+
 \frac{(\kappa^2-k^2)^2}{4k^2\kappa^2}\sinh^2\kappa a\\
 &=1+\frac{(\kappa^2+k^2)^2}{4k^2\kappa^2}\sinh^2\kappa a.
\end{align*}
使用$\kappa^2+k^2=2mV_0/\hbar^2$和$k^2\kappa^2=4m^2E(V_0-E)/\hbar^4$，得到
\[
 T=\left[1+\frac{V_0^2\sinh^2\kappa a}{4E(V_0-E)}\right]^{-1}.
\]
图\ref{fig:m4tunnela}的厚垒极限$\kappa a\gg1$给$T\simeq16E(V_0-E)\ee^{-2\kappa a}/V_0^2$。有限禁区内不能先删增长支；它与衰减支一起满足两端匹配，正如有限光隙中的受抑全反射。
\end{solution}

\begin{exercise}\label{ex:m4fieldemission}
金属表面外的电子势垒近似为$V(x)=W-eFx$，$0<x<W/(eF)$，其中$F>0$是外加电场大小。以表面内电子能量为零，求半经典穿透概率的主要指数，解释增大电场的作用。
\Diagram{fig:m4emissionq}{外场使表面势垒成为有限宽三角形}{\draw[axis](-.6,0)--(4.9,0)node[right]{$x$};\draw[axis](0,0)--(0,3.2)node[above]{$V$};\draw[curve](0,2.6)--(3.8,0);\node[left]at(0,2.6){$W$};\node[below]at(3.8,0){$W/(eF)$};\draw[->](-.5,.3)--(-.05,.3);}
\end{exercise}
\begin{solution}
图\ref{fig:m4emissionq}禁区局部虚动量大小$|p|=\sqrt{2m(W-eFx)}$。先计算禁区积分，取$q=W-eFx$，$\dd x=-\dd q/(eF)$，空间上下限对应$q=W,0$：
\begin{align*}
 \int_0^{W/(eF)}|p|\dd x
 &=-\frac{\sqrt{2m}}{eF}\int_W^0q^{1/2}\dd q\\
 &=\frac{\sqrt{2m}}{eF}\left[\frac23q^{3/2}\right]_0^W
 =\frac{2\sqrt{2m}W^{3/2}}{3eF}.
\end{align*}
WKB振幅的衰减因子为$\exp[-\int|p|\dd x/\hbar]$；概率需平方，故
\[
 \mathcal T\propto\exp\left[-\frac2\hbar\int_0^{W/(eF)}|p|\dd x\right]
 =\exp\left[-\frac{4\sqrt{2m}W^{3/2}}{3e\hbar F}\right].
\]
\Diagram{fig:m4emissiona}{电场增大使势垒变窄，指数抑制减弱}{\draw[axis](0,0)--(4.8,0)node[right]{$x$};\draw[axis](0,0)--(0,3.2)node[above]{$V$};\draw[curve](0,2.6)--(4,0);\draw[guide](0,2.6)--(2,0);\node[right]at(3,1.2){较弱场};\node[left]at(-.4,1.4){较强场};\draw[guide](-.3,1.4)--(.8,1.56);}
图\ref{fig:m4emissiona}中势垒高度仍为$W$，宽度却随$1/F$减小，所以穿透概率随电场迅速上升。若画$\ln\mathcal T$对$1/F$，领先指数给负斜率直线；实际电流还含电子供给与前因子，不能仅凭此式求出绝对电流。这个过程不要求入射光子能量超过功函数，和经典过垒光电逃逸属于不同模型。
\end{solution}

\begin{exercise}\label{ex:m4angularvariant}
光电角度模型改为只在向外半球产生电子，概率密度按$\cos\theta$分布，即$\dd P=C\cos\theta\dd\Omega$。法向势垒仍为$W$。求归一化常数和逃逸比例，与全空间各向同性模型比较。
\Diagram{fig:m4angularq}{方向权重偏向表面法线}{\draw[curve](-2,0)--(2,0);\draw[axis](0,0)--(0,2.6)node[above]{法线};\draw[->](0,0)--(.7,2.1);\draw[->](0,0)--(1.7,.8);\node[right]at(1.7,.8){$\cos\theta$权重};}
\end{exercise}
\begin{solution}
图\ref{fig:m4angularq}在向外半球归一化：
\[
 1=2\pi C\int_0^{\pi/2}\cos\theta\sin\theta\dd\theta
 =2\pi C\left[\frac{\sin^2\theta}{2}\right]_0^{\pi/2}
 =\pi C,
\]
故$C=1/\pi$。当$h\nu\geq W$时，可逸出角域为$0\leq\theta\leq\theta_c$，其中$\cos\theta_c=\sqrt{W/(h\nu)}$。在同一原函数中代入新的上限，
\[
 P_{\rm esc}=2\int_0^{\theta_c}\cos\theta\sin\theta\dd\theta
 =[\sin^2\theta]_0^{\theta_c}=1-\frac W{h\nu}.
\]
低于阈值时比例为零。
\Diagram{fig:m4angulara}{同一势垒，不同初始角分布产生不同逃逸曲线}{\draw[axis](0,0)--(4.6,0)node[right]{$h\nu/W$};\draw[axis](0,0)--(0,2.7)node[above]{$P_{\rm esc}$};\draw[curve,domain=1:4,samples=60]plot(\x,{2.3*(1-1/\x)});\draw[guide,domain=1:4,samples=60]plot(\x,{1.15*(1-sqrt(1/\x))});}
图\ref{fig:m4angulara}中此比例在高能趋1，而全空间各向同性模型趋$1/2$。势垒决定允许角域，产生机制决定各方向的权重。因此，相同能量阈值可以对应不同光电产额。
\end{solution}

\begin{exercise}\label{ex:m4radiativecool}
热容量$C$近似恒定的黑球表面积$A$，置于温度$T_0$的大黑腔。球温$T=T_0+\Delta T$，$|\Delta T|\ll T_0$。推导小温差弛豫时间，并说明面积和热容量怎样进入。
\Diagram{fig:m4coolq}{黑球与大腔壁双向辐射交换}{\draw[curve](0,0)circle(.8);\draw[curve](-2,-1.7)rectangle(2,1.7);\node at(0,0){$T$};\node[above]at(0,1.7){$T_0$};\draw[->](.9,.3)--(1.7,.3);\draw[->](-1.7,-.3)--(-.9,-.3);}
\end{exercise}
\begin{solution}
图\ref{fig:m4coolq}中，腔内平衡辐射各向同性，球吸收功率为$A\sigma T_0^4$，自身发射功率为$A\sigma T^4$。以整球热容量$C$列能量方程，得
\[
 C\frac{\dd T}{\dd t}=-A\sigma(T^4-T_0^4).
\]
写$T=T_0+\Delta T$，展开得$T^4-T_0^4=4T_0^3\Delta T+6T_0^2\Delta T^2+\cdots$。在$|\Delta T|/T_0\ll1$时保留线性项，
\[
 \frac{\dd\Delta T}{\dd t}=-\frac{\Delta T}{\tau},\qquad
 \tau=\frac{C}{4A\sigma T_0^3}.
\]
\Diagram{fig:m4coola}{小温差以指数形式衰减}{\draw[axis](0,0)--(4.5,0)node[right]{$t/\tau$};\draw[axis](0,0)--(0,2.6)node[above]{$\Delta T/\Delta T(0)$};\draw[curve,domain=0:4,samples=60]plot(\x,{2.2*exp(-\x)});}
当$\Delta T(0)\ne0$时分离变量，
\[
 \int_{\Delta T(0)}^{\Delta T(t)}\frac{\dd q}{q}
 =-\frac1\tau\int_0^t\dd s,
 \quad \ln\frac{|\Delta T(t)|}{|\Delta T(0)|}=-\frac t\tau.
\]
温差的符号不变，故$\Delta T(t)=\Delta T(0)\ee^{-t/\tau}$，得到图\ref{fig:m4coola}的曲线；初温恰等于腔温时，另有恒解$\Delta T=0$。初温低于腔温时，负温差也按同一式趋于零。量纲上$A\sigma T_0^3$为功率每温度，与$C$相除给时间。

对同材料均匀实心球，$C=\rho c_p(4\pi a^3/3)$、$A=4\pi a^2$，所以$\tau=\rho c_pa/(12\sigma T_0^3)$，半径越大弛豫越慢。该结论要求内部导热足够快，使整球可用单一温度描述。若环境温度可忽略，则原方程改为$\dot T=-(A\sigma/C)T^4$。分离变量，
\[
 \int_{T(0)}^{T(t)}T^{-4}\dd T=-\frac{A\sigma}{C}\int_0^t\dd s,
 \quad \left[-\frac1{3T^3}\right]_{T(0)}^{T(t)}=-\frac{A\sigma t}{C},
\]
得到$T^{-3}(t)=T^{-3}(0)+3A\sigma t/C$，为代数冷却。它采用不同的环境条件，不能沿用小温差指数律。
\end{solution}

\begin{exercise}\label{ex:m4spectralpeak}
求黑体“每对数频率间隔”的能量谱峰值方程，并说明它为什么不同于$u_\nu$的峰。
\Diagram{fig:m4logq}{相同对数宽度在高频处包含更大的线性频率区间}{\draw[axis](0,0)--(5.4,0)node[right]{$\nu$};\draw[very thick](.5,0)--(1,0);\draw[very thick](2,0)--(4,0);\node[above]at(.75,.5){$\nu\to2\nu$};\node[above]at(3,.5){$4\nu\to8\nu$};}
\end{exercise}
\begin{solution}
图\ref{fig:m4logq}比较等比例频率区间。因$\dd\nu=\nu\dd\ln\nu$，对数谱为$\nu u_\nu$，以$x=h\nu/(\kb T)$表示正比$f_4(x)=x^4/(\ee^x-1)$。在例\ref{ex:m4wien}已逐步求出的$f_j'$中取$j=4$，
\[
 f_4'(x)=\frac{x^3\ee^x}{(\ee^x-1)^2}
 [4(1-\ee^{-x})-x].
\]
除去边界根$0$，峰值方程为$4(1-\ee^{-x})=x$，非零正根$x\simeq3.92069$。
\Diagram{fig:m4loga}{对数谱多一个频率权重，峰向高频移动}{\draw[axis](0,0)--(5.3,0)node[right]{$h\nu/(\kb T)$};\draw[axis](0,0)--(0,2.8)node[above]{对数谱};\draw[curve,domain=.05:8,samples=80]plot({.6*\x},{.4*\x^4/(exp(\x)-1)});\draw[guide](2.3524,0)--(2.3524,2);}
图\ref{fig:m4loga}中的峰对应每相同倍频带的能量贡献，而$u_\nu$比较每相同赫兹的能量。它也等于每对数波长间隔能量谱的峰所对应的同一物理频率，因为$|\dd\ln\lambda|=|\dd\ln\nu|$；线性波长谱则还有不同雅可比。
\end{solution}

\begin{exercise}\label{ex:m4welltransition}
无限深方阱宽度缓慢从$L$变为$2L$，粒子始终留在同一量子数$n$的定态分支。求能量与典型动量的变化，并用作用量解释“缓慢”的意义。
\Diagram{fig:m4expandq}{缓慢移动硬壁使阱宽加倍}{\draw[curve](0,2)--(0,0)--(2,0)--(2,2);\draw[guide](4,0)--(4,2);\draw[->](2.2,1)--(3.8,1);\node[below]at(2,-.4){$L\to2L$};}
\end{exercise}
\begin{solution}
图\ref{fig:m4expandq}的定态能量$E_n=n^2h^2/(8mL^2)$，所以阱宽加倍，能量减为$1/4$，动量大小$p_n=nh/(2L)$减为$1/2$。定态内平均动量为0，但$\langle p^2\rangle^{1/2}=p_n$，不要把典型动量大小等同于平均动量。
\Diagram{fig:m4expanda}{相轨道长宽反比变化，面积$2pL=nh$保持}{\draw[axis](0,0)--(5,0)node[right]{$x$};\draw[axis](0,-1.8)--(0,2)node[above]{$p$};\draw[curve](0,-1.3)rectangle(2,1.3);\draw[guide](0,-.65)rectangle(4,.65);}
图\ref{fig:m4expanda}说明缓变过程中作用量$J=2pL$近似保持，量子上要求壁运动足够慢，避免显著跃迁到别的$n$。若壁突然移动，原波函数不再是新阱单一定态，需要展开为新本征态之和，本题能量比例就不能直接使用。
\end{solution}

\begin{exercise}\label{ex:m4oscforce}
谐振子势中加入恒定力$F$，势为$m\omega^2x^2/2-Fx$。求能级、平衡位置和能级间隔，说明零点项是否改变。
\Diagram{fig:m4forceq}{恒力使抛物势横移并下移}{\draw[axis](-2.2,0)--(3.4,0)node[right]{$x$};\draw[axis](0,-1)--(0,3)node[above]{$V$};\draw[guide,domain=-1.8:1.8,samples=50]plot(\x,{.7*\x*\x});\draw[curve,domain=-1:2.8,samples=50]plot(\x,{.7*(\x-.9)^2-.567});}
\end{exercise}
\begin{solution}
逐项求导得$V'(x)=m\omega^2x-F$，令其为零，平衡位置为$x_0=F/(m\omega^2)$；再求导得$V''=m\omega^2>0$，故这是稳定平衡。配方得$V=\tfrac12m\omega^2(x-x_0)^2-F^2/(2m\omega^2)$。图\ref{fig:m4forceq}中的势曲率不变，运动频率仍为$\omega$，能级为$E_n=\hbar\omega(n+1/2)-F^2/(2m\omega^2)$。
\Diagram{fig:m4forcea}{所有能级整体平移，间隔$\hbar\omega$保持}{\foreach\y in{0,.7,1.4,2.1}{\draw[curve](-2,\y)--(-.3,\y);\draw[curve](1,{\y-.5})--(2.7,{\y-.5});}\draw[->](0,1.4)--(.7,.9);\node[above]at(1.8,2.3){有恒力};}
图\ref{fig:m4forcea}中零点能相对势阱最低点仍为$\hbar\omega/2$，但相对原来的能量零点发生共同平移。边界仍是两个平滑转点，故作用量偏移$1/2$不变。只看总能量可能出现负值，并不表示振动激发能或频率为负。
\end{solution}

\begin{exercise}\label{ex:m4hydrogenscaling}
氢样离子核电荷由$Z$变为$2Z$，量子数保持$n$。分别讨论非相对论玻尔半径、速度、束缚能和相邻能级频率的缩放，指出何时这种缩放失效。
\Diagram{fig:m4Zq}{同一$n$下增大核电荷，轨道缩小}{\draw[curve](-1.5,0)circle(1.1);\fill(-1.5,0)circle(2pt);\draw[curve](2,0)circle(.55);\fill(2,0)circle(2pt);\node[below]at(-1.5,-1.3){$Z$};\node[below]at(2,-1.3){$2Z$};}
\end{exercise}
\begin{solution}
图\ref{fig:m4Zq}依$r_n\propto1/Z$、$v_n\propto Z$、$E_n\propto-Z^2$，半径减半、速率加倍、束缚能绝对值乘4。固定两量子数的谱线差同乘4，频率也乘4。
\Diagram{fig:m4Za}{速度尺度$Z\alpha/n$决定非相对论近似的适用性}{\draw[axis](0,0)--(4.5,0)node[right]{$Z\alpha/n$};\draw[axis](0,0)--(0,2.7)node[above]{$v/c$};\draw[curve](0,0)--(3.7,2.3);\node[right]at(3.7,2.3){近1时须改模型};}
图\ref{fig:m4Za}中$Z\alpha/n\ll1$是非相对论条件。增大$Z$会增强相对论修正；当参数逼近1，圆轨道半经典模型本身也失去常规适用性，有限核尺寸、自旋和完整量子动力学不可忽略。缩放律是在固定模型条件下的结果，不是任意核电荷都可无限外推。
\end{solution}

\begin{exercise}\label{ex:m4yield}
采用例\ref{ex:m4escape}的角度筛选模型，在固定吸收功率$P$下改变频率。求零电压光电流最大的光子能量，以及最大电流。
\Diagram{fig:m4yieldq}{光子能量增加时每秒光子数减少，逃逸锥却扩大}{\draw[axis](0,0)--(0,2.7)node[above]{法线};\draw[curve](0,0)--(.8,2);\draw[curve](0,0)--(1.8,2);\node[right]at(1.8,2){较高$h\nu$};\node[left]at(-.1,1){$P$固定};}
\end{exercise}
\begin{solution}
令$x=h\nu/W\geq1$，图\ref{fig:m4yieldq}的逃逸概率为$(1-x^{-1/2})/2$，所以$I=eP( x^{-1}-x^{-3/2})/(2W)$。求导得$\dd I/\dd x=(eP/2W)(-x^{-2}+\tfrac32x^{-5/2})$，零点为$\sqrt x=3/2$。
\Diagram{fig:m4yielda}{电流在$x=9/4$达到极大}{\draw[axis](0,0)--(5,0)node[right]{$x=h\nu/W$};\draw[axis](0,0)--(0,2.6)node[above]{相对电流};\draw[curve,domain=1:7,samples=90]plot({.6*\x},{14*(1/\x-1/(\x^1.5))});\draw[guide](1.35,0)--(1.35,2.075);\node[below]at(1.35,0){$9/4$};}
在$x=1$电流为0，$x\to\infty$也趋于0，故图\ref{fig:m4yielda}的驻点是全局极大。代入得$I_{\max}=2eP/(27W)$、$h\nu=9W/4$。若只看逃逸概率而忽略光子数$P/(h\nu)$，就会错误地认为频率越高电流必越大。
\end{solution}

\begin{exercise}\label{ex:m4cavitypressure}
各向同性平衡辐射能量密度为$u$，理想反射壁法线为$z$。由光子动量通量求辐射压强；再求一束垂直单向光被完全吸收、完全反射的压强。
\Diagram{fig:m4pressureq}{斜向光子反射时法向动量改变$2p\cos\theta$}{\draw[curve](-2,0)--(2,0);\draw[->](-1.6,1.8)--(0,0);\draw[->](0,0)--(1.6,1.8);\draw[guide](0,0)--(0,2.4);\node at(.45,1.3){$\theta$};}
\end{exercise}
\begin{solution}
图\ref{fig:m4pressureq}中一个光子的法向动量改变$2E\cos\theta/c$。单位时间单位面积入射光子的能量还带通量投影$c\cos\theta$，因此
\[
 p=\frac{2u}{4\pi}\int_0^{2\pi}\dd\varphi\int_0^{\pi/2}\cos^2\theta\sin\theta\dd\theta
 =u\left[-\frac{\cos^3\theta}{3}\right]_0^{\pi/2}=\frac u3.
\]
\Diagram{fig:m4pressurea}{单向垂直光束的吸收和反射动量通量}{\draw[curve](0,-.5)--(0,1.2);\draw[->](-2,.8)--(-.1,.8)node[midway,above]{$I/c$};\draw[->](-.1,.2)--(-2,.2)node[midway,below]{反射多一份};\node[right]at(.2,.6){壁};}
图\ref{fig:m4pressurea}的单向光束不需要半球积分，完全吸收时$p=I/c$，完全反射时$p=2I/c$。这里$I$是垂直于光束的功率面密度。腔辐射的$u/3$与孔口通量$cu/4$来自不同角度矩：前者有$\cos^2\theta$，后者有$\cos\theta$。
\end{solution}

\begin{exercise}\label{ex:m4equilibrium}
半径$R$的恒星近似黑体温度$T_s$，距离$r\gg R$处有半径$a$、表面吸收和发射均为黑体的小球。忽略其他热源，求球均匀温度。若只在受光半球均匀发热、背面不辐射，结果如何？
\Diagram{fig:m4equilibriumq}{恒星辐射照到小球的截面积$\pi a^2$}{\draw[curve](0,0)circle(.8);\node at(0,0){$T_s$};\draw[curve](5,0)circle(.35);\foreach\y in{-.2,.2}{\draw[->](1,\y)--(4.6,\y);}\node[below]at(2.8,-.5){$r$};\node[right]at(5.4,0){$a$};}
\end{exercise}
\begin{solution}
恒星总功率$4\pi R^2\sigma T_s^4$，在半径$r$球面上的通量$F=\sigma T_s^4R^2/r^2$。图\ref{fig:m4equilibriumq}的小球吸收面积是迎光截面$\pi a^2$，吸收功率$\pi a^2F$。
\Diagram{fig:m4equilibriuma}{吸收用投影面积，发射用实际表面积}{\draw[curve](0,0)circle(1);\draw[guide](0,-1)--(0,1);\draw[->](1,0)--(1.7,0);\draw[->](0,1)--(0,1.7);\draw[->](-1,0)--(-1.7,0);\draw[->](0,-1)--(0,-1.7);\node[right]at(2,0){全表面$4\pi a^2$};}
均匀温度时，图\ref{fig:m4equilibriuma}全表面发射$4\pi a^2\sigma T^4$。平衡给$T=T_s\sqrt{R/(2r)}$。若只有半球辐射，发射面积$2\pi a^2$，则$T=2^{1/4}T_s\sqrt{R/(2r)}$。该变式是规定的温度分布模型，真实表面是否均温取决于导热、自转等过程。
\end{solution}

\begin{exercise}\label{ex:m4reducedmass}
把玻尔模型的固定核改为质量$M$可动的原子核。求相对坐标的半径、能级及相对固定核结果的修正。
\Diagram{fig:m4reducedq}{电子和核绕共同质心转动}{\draw[curve](0,0)circle(1.7);\draw[curve](0,0)circle(.4);\fill(1.7,0)circle(2pt)node[right]{$m$};\fill(-.4,0)circle(3pt)node[above left]{$M$};\fill(0,0)circle(1pt);\draw[guide](0,0)--(.6,-.65)node[right]{质心};}
\end{exercise}
\begin{solution}
图\ref{fig:m4reducedq}中记电子、原子核位置为$\vct r_e,\vct r_N$，定义质心$\vct R=(m\vct r_e+M\vct r_N)/(m+M)$和相对坐标$\vct r=\vct r_e-\vct r_N$。反解为
\[
 \vct r_e=\vct R+\frac M{M+m}\vct r,\quad
 \vct r_N=\vct R-\frac m{M+m}\vct r.
\]
代入动能，交叉项相消，得到
\[
 K=\tfrac12(M+m)\dot{\vct R}^{\,2}
 +\tfrac12\mu\dot{\vct r}^{\,2},\qquad \mu=\frac{mM}{m+M}.
\]
库仑势仅依赖$r$，因此相对运动等价于质量$\mu$在固定库仑中心中运动。
\Diagram{fig:m4reduceda}{核、电子到质心距离按反质量比分配}{\draw[axis](-1.7,0)--(3.9,0);\fill(-1,0)circle(3pt)node[above]{$M$};\fill(0,0)circle(1.5pt)node[above]{$C$};\fill(3,0)circle(2pt)node[above]{$m$};\draw[<->](-1,-.5)--(0,-.5)node[midway,below]{$r_N$};\draw[<->](0,-.5)--(3,-.5)node[midway,below]{$r_e$};}
在质心系对相对角动量取$\mu vr=n\hbar$，向心方程为$\mu v^2/r=Ze^2/(4\pi\varepsilon_0r^2)$。两式相除先得$v=Z\alpha c/n$，再得
\[
 r_n=\frac{4\pi\varepsilon_0\hbar^2n^2}{\mu Ze^2},\qquad
 E_n=-\frac{\mu c^2(Z\alpha)^2}{2n^2}.
\]
图\ref{fig:m4reduceda}中的电子、核到质心距离分别为$Mr_n/(M+m)$、$mr_n/(M+m)$。相对固定核结果，能级绝对值乘$\mu/m=1/(1+m/M)$，核间距乘其倒数。若$M\gg m$，能级相对修正为$-m/M+O(m^2/M^2)$，相应跃迁频率也降低同一比例。核可动导致的约化质量修正改变内部能级；整个原子发光时的反冲则还需要在这些能级之上作总四动量守恒。
\end{solution}

\begin{exercise}\label{ex:m4powerwell}
粒子在一维势$V(x)=A|x|^s$中运动，$A>0,s>0$。用两平滑转点条件给出能级对量子数的幂次关系，并比较$s=2$与$s\to\infty$的极限。
\Diagram{fig:m4powerq}{幂次束缚势的对称转点}{\draw[axis](-2,0)--(2.3,0)node[right]{$x$};\draw[axis](0,0)--(0,3)node[above]{$V$};\draw[curve,domain=-1.7:1.7,samples=70]plot(\x,{.5*abs(\x)^3});\draw[guide](-1.5,1.6875)--(1.5,1.6875);\node[right]at(1.5,1.6875){$E$};}
\end{exercise}
\begin{solution}
图\ref{fig:m4powerq}中转点为$\pm x_t$，$x_t=(E/A)^{1/s}$。利用势的偶对称，作用量写为
\[
 J=4\int_0^{x_t}\sqrt{2m(E-Ax^s)}\dd x.
\]
令$y=x/x_t$，有$\dd x=x_t\dd y$、$Ax^s=Ey^s$，所以
\[
 J=4\sqrt{2m}A^{-1/s}E^{(s+2)/(2s)}C_s,\quad
 C_s=\int_0^1\sqrt{1-y^s}\dd y.
\]
\Diagram{fig:m4powera}{作用量积分换成固定区间的无量纲面积}{\begin{scope}[yscale=.45]\draw[axis](0,0)--(3.9,0)node[right]{$y$};\draw[axis](0,0)--(0,2.5)node[above]{$\sqrt{1-y^s}$};\draw[curve,domain=0:1,samples=70]plot({3*\x},{2*sqrt(1-\x^3)});\node[below]at(3,0){$1$};\end{scope}}
图\ref{fig:m4powera}的无量纲面积与$E$无关。两软转点条件给出
\[
 E_n\simeq\left[\frac{hA^{1/s}}{4\sqrt{2m}C_s}
 (n+1/2)\right]^{2s/(s+2)}.
\]
因此$s<2$时高能级间距减小，$s=2$时等间隔，$s>2$时增大。取$s=2$，$C_2=\pi/4$；再令$A=m\omega^2/2$，系数化为$\hbar\omega$，恢复谐振子结果。

讨论$s\to\infty$时应先固定长度尺度，写$V=V_*|x/L_*|^s$。对有限$E$，转点趋$L_*$、$C_s\to1$，作用量趋$4L_*\sqrt{2mE}$，与宽$2L_*$的方阱相同，能量的高量子数标度为$n^2$。但端点在极限中变为硬壁，量子数偏移须重新匹配，不能把软转点的$1/2$保留为硬壁最低能级的准确偏移。原点在$s\leq1$时不光滑，转点仍平滑；相关半经典式主要用于高能级，低能级须另解边界问题。
\end{solution}

\begin{exercise}\label{ex:m4gravityscale}
定义重力反弹模型的特征长度$\ell=(\hbar^2/(2m^2g))^{1/3}$。把其薛定谔方程化为无量纲形式，说明精确能量由何种边界条件确定，并比较半经典结果。
\Diagram{fig:m4gravityq}{硬壁处波函数为零，高处波函数衰减}{\draw[axis](-.3,0)--(4.7,0)node[right]{$x$};\node[above]at(0,1.6){$V$};\draw[very thick](0,0)--(0,1.4);\draw[curve](0,0)--(4,1.3);\draw[guide](0,.9)--(4.3,.9);\node[right]at(4.3,.9){$E$};}
\end{exercise}
\begin{solution}
图\ref{fig:m4gravityq}对应$-\hbar^2\psi''/(2m)+mgx\psi=E\psi$。令$\xi=x/\ell$、$\epsilon=E/(mg\ell)$，则链式法则给$\dd\psi/\dd x=\ell^{-1}\dd\psi/\dd\xi$、$\dd^2\psi/\dd x^2=\ell^{-2}\dd^2\psi/\dd\xi^2$。代入并除以$mg\ell$，用$\hbar^2/(2m\ell^2)=mg\ell$，化为
\[
 -\frac{\dd^2\psi}{\dd\xi^2}+\xi\psi=\epsilon\psi.
\]
\Diagram{fig:m4gravitya}{无量纲重力反弹基态在硬壁为零，在高处衰减}{\begin{scope}[yscale=.65]\draw[axis](0,0)--(7.5,0)node[right]{$\xi$};\draw[axis](0,-.4)--(0,2.1)node[above]{$\psi$};\draw[curve]plot coordinates {(0.00000,-0.0000000) (0.06422,0.1348816) (0.12844,0.2685052) (0.19266,0.3996896) (0.25688,0.5273458) (0.32110,0.6504842) (0.38532,0.7682206) (0.44954,0.8797792) (0.51376,0.9844945) (0.57798,1.0818115) (0.64220,1.1712844) (0.70642,1.2525734) (0.77064,1.3254415) (0.83486,1.3897488) (0.89908,1.4454471) (0.96330,1.4925734) (1.02752,1.5312426) (1.09174,1.5616405) (1.15596,1.5840157) (1.22018,1.5986721) (1.28440,1.6059610) (1.34862,1.6062736) (1.41284,1.6000332) (1.47706,1.5876887) (1.54128,1.5697070) (1.60550,1.5465674) (1.66972,1.5187550) (1.73394,1.4867556) (1.79817,1.4510510) (1.86239,1.4121141) (1.92661,1.3704052) (1.99083,1.3263691) (2.05505,1.2804315) (2.11927,1.2329971) (2.18349,1.1844479) (2.24771,1.1351413) (2.31193,1.0854092) (2.37615,1.0355578) (2.44037,0.9858666) (2.50459,0.9365891) (2.56881,0.8879525) (2.63303,0.8401586) (2.69725,0.7933847) (2.76147,0.7477840) (2.82569,0.7034872) (2.88991,0.6606034) (2.95413,0.6192213) (3.01835,0.5794110) (3.08257,0.5412249) (3.14679,0.5046992) (3.21101,0.4698558) (3.27523,0.4367028) (3.33945,0.4052368) (3.40367,0.3754436) (3.46789,0.3472996) (3.53211,0.3207732) (3.59633,0.2958260) (3.66055,0.2724135) (3.72477,0.2504864) (3.78899,0.2299915) (3.85321,0.2108727) (3.91743,0.1930715) (3.98165,0.1765277) (4.04587,0.1611805) (4.11009,0.1469685) (4.17431,0.1338306) (4.23853,0.1217063) (4.30275,0.1105360) (4.36697,0.1002615) (4.43119,0.0908260) (4.49541,0.0821748) (4.55963,0.0742549) (4.62385,0.0670155) (4.68807,0.0604080) (4.75229,0.0543862) (4.81651,0.0489061) (4.88073,0.0439261) (4.94495,0.0394069) (5.00917,0.0353115) (5.07339,0.0316054) (5.13761,0.0282559) (5.20183,0.0252329) (5.26606,0.0225080) (5.33028,0.0200552) (5.39450,0.0178499) (5.45872,0.0158698) (5.52294,0.0140941) (5.58716,0.0125036) (5.65138,0.0110808) (5.71560,0.0098096) (5.77982,0.0086750) (5.84404,0.0076638) (5.90826,0.0067634) (5.97248,0.0059627) (6.03670,0.0052514) (6.10092,0.0046203) (6.16514,0.0040610) (6.22936,0.0035658) (6.29358,0.0031280) (6.35780,0.0027412) (6.42202,0.0023999) (6.48624,0.0020990) (6.55046,0.0018341) (6.61468,0.0016011) (6.67890,0.0013964) (6.74312,0.0012167) (6.80734,0.0010592) (6.87156,0.0009211) (6.93578,0.0008004) (7.00000,0.0006948)};\node[below left]at(0,0){$\psi(0)=0$};\node[above]at(4.6,.6){$\operatorname{Ai}(\xi-\epsilon_0)$};\end{scope}}
令$z=\xi-\epsilon$，方程为$\psi''=z\psi$，通解为$A\operatorname{Ai}(z)+B\operatorname{Bi}(z)$。正无穷处$\operatorname{Bi}$增长，归一化要求$B=0$；$A\ne0$再由归一化确定。图\ref{fig:m4gravitya}的硬壁$x=0$对应$z=-\epsilon$，要求$\operatorname{Ai}(-\epsilon)=0$，所以精确能量由其负零点确定。半经典近似给$\epsilon_n\simeq[\tfrac{3\pi}2(n+\tfrac34)]^{2/3}$。无量纲化把质量和重力全部收进尺度$mg\ell$，可用于比较不同粒子而不重复解微分方程。
\end{solution}

